% ===================================================================
% chapters/04_derivations_electrostatics.tex  (writer D1)
% Physical model I: electrostatics, band structure, density of states and charges
% ASCII only.
% ===================================================================
\chapter{Physical model I: electrostatics, band structure, density of states and charges}
\label{ch:derivations-es}

This chapter documents, one input at a time, every parameter that determines the electrostatics of the
simulated IWO thin-film transistor: the gate stack, the dielectric constant and band structure of the
channel, the thickness-dependent conduction-band shift attributed to quantum confinement, the effective
masses and effective densities of states, the localized-state (trap) spectrum, and the fixed and
interface charges. It closes with three derivations that relate these inputs to the measured quantities:
the threshold charge budget, the subthreshold-slope relations (including the trap-density inversion
used in the analysis), and the Debye and depletion length scales that decide whether a film is fully depleted.
For each electrostatic input the chapter states the number that enters the ATLAS deck, its origin, the way in
which it could be derived or measured, the strength of its control over the simulated transfer curve, and the
confidence that it deserves.

Every section has the same internal structure: definition and role, values per thickness, source,
derivation, ATLAS implementation, sensitivity evidence and status. Status labels follow the provenance
vocabulary (\prov{MEASURED\_TARGET\_DEVICE}, \prov{LITERATURE\_IWO}, \prov{MEASURED\_RELATED\_MATERIAL},
\prov{DFT\_PURE\_IN2O3\_PROXY}, \prov{FITTED}, \prov{ASSUMED}, \prov{CALCULATED}); evidence classes use
\evDATA, \evNUM, \evCAL, \evOOS, \evPOST, \evPRED, \evSURR, \evTHEORY, \evLIT{} as defined in the nomenclature.
The configuration values are those of \file{config/iwo_material_model.yaml}. The analysis reports of
2026-09-25 are abbreviated as S1 (\file{analysis_2026-09-25/S1_electrostatics_6p3nm.md}), S3
(\file{S3_quantum_confinement.md}), S4 (\file{S4_defects_traps.md}), S6
(\file{S6_tcad_numerics_predictions.md}), the evidence brief (\file{EVIDENCE_BRIEF.md}), REVIEW and
SYNTHESIS, all in \file{analysis_2026-09-25/}. Sensitivity numbers at 2~nm are one-at-a-time
perturbations of the calibrated deck \run{0012} (DOS 96/48) or, for the isolation tests, of \run{0038}
(DOS 384/192); they are summarised graphically in \cref{fig:param_control_2nm,fig:param_control_delta}
and discussed as a whole in \cref{sec:cal-param-control}. Three standing corrections of SYNTHESIS
section~K apply throughout: \run{0027} (contact overlap) is malformed and retracted and is never used
here; $\SSmin$ is not a reliable metric (\cref{sec:der-metrics}), so slopes are quoted at fixed current; and the confinement shift
is not identified by the transfer curves, because it is degenerate with a fixed charge
(\cref{sec:par-dec,sec:par-qf}).

Throughout, $q$ is the elementary charge, $\kT = \SI{25.85}{meV}$ ($T = \SI{300}{K}$ assumed; the
measurement temperature was not recorded), $\tIWO$ (or $t$) the nominal channel thickness, and energies of
localized states are referenced to the local conduction-band edge $\Ec$ unless stated otherwise.

% -------------------------------------------------------------------
\section{Gate stack: thicknesses, permittivities, oxide capacitance and EOT}
\label{sec:par-stack}

\paragraph{Definition and role.}
The bottom gate is a TiN electrode (\SI{50}{nm}) under a bilayer dielectric of \HfO{} (\SI{15}{nm}) and
\AlO{} (\SI{2}{nm}); the IWO channel sits on the \AlO{}, Pd source/drain contacts (\SI{70}{nm}) are on top,
and the channel is exposed to air (no passivation); $W/L = 290/20~\mu$m, $\Vd = \SI{0.7}{V}$
(\cref{fig:stack_schematic}, \cref{sec:inp-geometry}). The gate-stack capacitance $\Cox$ converts every sheet
charge into a gate-voltage shift and is therefore the most important electrostatic constant of the
model; the ratio $q/\Cox$ converts charge into threshold voltage in every section below.

\paragraph{Values used.}
Thickness-independent (the same stack for all films): $t_{\mathrm{HfO_2}} = \SI{15}{nm}$,
$\varepsilon_{\mathrm{HfO_2}} = 19.57$, $t_{\mathrm{Al_2O_3}} = \SI{2}{nm}$,
$\varepsilon_{\mathrm{Al_2O_3}} = 9.0$; hence $\EOT = \SI{3.86}{nm}$,
$\Cox = \SI{8.955e-7}{\farad\per\centi\metre\squared}$ and $q/\Cox = \SI{0.179}{V}$ per
$10^{12}\cmsq$ (evidence brief).

\paragraph{Source.}
Layer thicknesses and device geometry: the device schematic \citep{DeviceSchematic}. The \HfO{}
permittivity 19.57 is the C--V value of the target process in the supporting information of the target
paper \citep[SI Fig.~S1b]{JanuarSI2026}; the SI is not bundled and the value was taken over from the Astra
audit (\file{digest/used_items.csv}). The \AlO{} permittivity 9.0 is \prov{ASSUMED} (not measured on the
target process).

\paragraph{Derivation.}
The two dielectrics are capacitors in series,
\begin{equation}
  \frac{1}{\Cox} = \frac{t_{\mathrm{HfO_2}}}{\varepsilon_0\varepsilon_{\mathrm{HfO_2}}}
                 + \frac{t_{\mathrm{Al_2O_3}}}{\varepsilon_0\varepsilon_{\mathrm{Al_2O_3}}},
  \qquad
  \EOT = 3.9\,\varepsilon_0/\Cox
       = t_{\mathrm{HfO_2}}\frac{3.9}{\varepsilon_{\mathrm{HfO_2}}} + t_{\mathrm{Al_2O_3}}\frac{3.9}{\varepsilon_{\mathrm{Al_2O_3}}}.
  \label{eq:d1-cox}
\end{equation}
With the values above, $\EOT = 15\times 3.9/19.57 + 2\times 3.9/9.0 = 2.989 + 0.867 = \SI{3.856}{nm}$ and
$\Cox = 3.9\times\SI{8.854e-14}{F/cm}/\SI{3.856e-7}{cm} = \SI{8.955e-7}{\farad\per\centi\metre\squared}$.
A sheet charge $\Delta Q$ (in cm$^{-2}$) placed at the channel/dielectric interface shifts the gate voltage by
\begin{equation}
  \Delta V = \frac{q\,\Delta Q}{\Cox} = \SI{0.1789}{V}\times\frac{\Delta Q}{10^{12}\cmsq}.
  \label{eq:d1-qcox}
\end{equation}
The dominant uncertainty is $\varepsilon_{\mathrm{Al_2O_3}}$: evaluating \cref{eq:d1-cox} with 8.0 gives
$\EOT = \SI{3.96}{nm}$ (+2.8\,\%) and with 7.0 gives \SI{4.10}{nm} (+6.4\,\%). REVIEW quotes ``about 3\,\%''
for the range 7--9; the calculation above shows that this holds for the upper half of the range, and that
the lower end reaches about 6\,\%. Every charge derived through $q/\Cox$ in this report (e.g.\ $\Qf$) scales with it.
An accumulation C--V measurement on a MIM capacitor of the same stack (TiN/\HfO/\AlO/metal)
would fix $\Cox$ directly, independently of both permittivities.

\paragraph{ATLAS implementation.}
The \HfO{} and \AlO{} layers are insulator REGIONs whose permittivity is set with the PERMITTIVITY parameter of
the MATERIAL statement \citep[p.~1340]{AtlasManual}; the TiN gate is an electrode with a work function
(\cref{sec:par-gatewf}). The full statement list is reproduced in \cref{ch:app-decks}.

\paragraph{Sensitivity evidence.}
No run varies $\Cox$; its effect enters analytically through \cref{eq:d1-qcox}. The rigid-shift property of
$\Cox$-weighted charges was verified in ATLAS: a change of $\Qf$ reproduces $\Delta\Vth = -q\Delta\Qf/\Cox$
to \SI{1}{mV} (runs \run{0001}$\rightarrow$\run{0003} and \run{0005}$\rightarrow$\run{0007}, evidence brief).

\paragraph{Status.}
Thicknesses \prov{MEASURED\_TARGET\_DEVICE} (nominal); $\varepsilon_{\mathrm{HfO_2}}$
\prov{MEASURED\_TARGET\_DEVICE}; $\varepsilon_{\mathrm{Al_2O_3}}$ \prov{ASSUMED}; $\Cox$, $\EOT$
\prov{CALCULATED}. Caveat: the thickness-measurement method is not reported, and 3--6\,\% of $\Cox$
uncertainty propagates into every fitted charge.

% -------------------------------------------------------------------
\section{Relative permittivity of the IWO channel}
\label{sec:par-eps-iwo}

\paragraph{Definition and role.}
$\epsIWO$ is the static relative permittivity of the channel. It enters Poisson's equation inside the film
and therefore only the \emph{internal} potential drop across the film: the back-surface lever arm
$t/\varepsilon_s$ (\cref{sec:par-backq}), the donor term $q\Nd t^2/2\varepsilon_s$ (\cref{sec:par-nd}), the
Debye and trap-screening lengths (\cref{sec:der-depletion}) and the image-charge correction to
confinement (\cref{sec:der-confinement}).

\paragraph{Values used.}
$\epsIWO = 9.3$ for all thicknesses; no thickness dependence is imposed (no evidence for one).

\paragraph{Source.}
C--V of IWO of the target process, $9.30\pm0.04$ \citep[SI Fig.~S1b]{JanuarSI2026} (SI not bundled; value
recorded by the Astra audit). The same SI gives 9.03 for \InO{} of the target process. Alternatives: the
static (DC) permittivity of single-crystal \InO{}, 10.55 \citep[p.~225102-12]{Stokey2021}, and 8.9 for
\InO{} as quoted by \citet[p.~501]{Si2021} (their source, Hamberg and Granqvist, concerns evaporated Sn-doped \InO{}
films; reference list p.~506).

\paragraph{Derivation (how it is measured).}
For a MIM or MIS structure with a film of thickness $t$, $\varepsilon_r = C t/(\varepsilon_0 A)$ in
depletion/insulating conditions; for a semiconducting film it is extracted from the series combination
with the known dielectric. A THz/Lyddane--Sachs--Teller route gives the crystal value,
$\varepsilon_{DC} = \varepsilon_\infty\prod(\omega_{LO}/\omega_{TO})^2$, used by \citet[p.~225102-12]{Stokey2021}.
The model uses the device-process C--V value because it contains the actual film microstructure.

\paragraph{ATLAS implementation.}
MATERIAL \ldots\ PERMITTIVITY=9.3 on the IWO region(s) \citep[p.~1340]{AtlasManual}.

\paragraph{Sensitivity evidence.}
\run{0026} ($\epsIWO = 10.55$ vs \run{0012}, 2~nm): $\Delta\Vthcc = -0.001$~V, fixed-current SS
($10^{-10}$--$10^{-9}\Aum$) $-0.4\mVdec$, $\Delta\Ion = +0.7$\,\% (\file{figures/MANIFEST.md}, param\_control\_delta).
At 2~nm the permittivity is electrostatically irrelevant (SYNTHESIS: ``donors and permittivity change the 2~nm
result by $<2$\,\%'', runs \run{0022}, \run{0026}). The internal term grows with $t$ (it is $\propto t/\varepsilon_s$),
so the insensitivity is not established at 13.2~nm, where no permittivity run exists.

\paragraph{Status.}
\prov{MEASURED\_TARGET\_DEVICE} (via an unbundled SI; page-level verification of the SI not possible here);
\evNUM{} for the insensitivity at 2~nm.

% -------------------------------------------------------------------
\section{Band gap}
\label{sec:par-eg}

\paragraph{Definition and role.}
$\Eg(t) = \Eg^{\mathrm{bulk}} + \dEg(t)$ is the fundamental gap of the channel. In the electrons-only
drift--diffusion model holes are not solved, so $\Eg$ enters only (i) the intrinsic density
$n_i = \sqrt{\Nc\Nv}\exp(-\Eg/2\kT)$, which is negligible, and (ii) the energy range over which the tail and
Gaussian defect distributions are integrated (\cref{sec:par-tail,sec:par-deep}). The thickness dependence is
carried entirely by the conduction band (\cref{sec:par-dec}).

\paragraph{Values used.}
$\Eg^{\mathrm{bulk}} = \SI{3.05}{eV}$ ($\pm\SI{0.2}{eV}$). With $\dEg = \dEc$ from \cref{eq:dec-law}:
$\Eg = 3.403$ / 3.122 / 3.076~eV at 2.0 / 6.3 / 13.2~nm (evaluated from the configuration law).

\paragraph{Source.}
Sputtered amorphous IWO of a different device, \citet[Sec.~3]{Fan2021} (not bundled; value taken over
from the Astra audit, \file{digest/used_items.csv}). PBE absolute gaps (0.94~eV bulk \InO{}; 1.27 and 1.88~eV
for 1.98 and 0.95~nm slabs, \citealp[p.~21542]{Lin2022}) are \emph{not} used because semi-local DFT
underestimates absolute gaps; only slab-minus-bulk differences enter (\cref{sec:par-dec}). The partition
$\dEc/\dEg = 1$ follows the observation that the valence-band edge is almost unchanged under confinement in
the \InO/\AlO{} slabs \citep[p.~504]{Si2021}.

\paragraph{Derivation (how it could be measured).}
Optical absorption (Tauc analysis), spectroscopic ellipsometry or REELS on films of each thickness would
give $\Eg(t)$; S3 lists the expected signature: $\Eg(2)-\Eg(13.2) \approx 0.2$--$0.45$~eV if the film
confinement is real, $<0.02$~eV if not (S3 falsification table).

\paragraph{ATLAS implementation.}
MATERIAL EG300 on each IWO region \citep[p.~1338]{AtlasManual}, set per thickness by the deck builder.

\paragraph{Sensitivity evidence.}
No run varies $\Eg$ alone. Since holes are not solved and $n_i \lesssim 10^{-7}\cmcu$ (evaluated from the
formula with the values of \cref{sec:par-nc,sec:par-nv}), $\Eg$ has no first-order effect on the transfer
curve; its only role is the integration range of the defect distributions, whose integrands are negligible
at mid-gap.

\paragraph{Status.}
$\Eg^{\mathrm{bulk}}$ \prov{LITERATURE\_IWO} (different device, not bundled); $\dEg(t)$
\prov{DFT\_PURE\_IN2O3\_PROXY}; the transfer curves are insensitive to it (\evTHEORY).

% -------------------------------------------------------------------
\section{Electron affinity}
\label{sec:par-chi}

\paragraph{Definition and role.}
$\chiIWO(t) = \chi^{\mathrm{bulk}} - \dEc(t)$ places the conduction-band edge relative to the vacuum level and
thereby relative to the gate work function $\phi_M$. The combination $\phi_M - \chiIWO$ is the flat-band
offset of the gate with respect to the channel conduction band; it enters the threshold voltage one-to-one
(\cref{sec:der-vth-budget}).

\paragraph{Values used.}
$\chi^{\mathrm{bulk}} = \SI{4.3}{eV}$ ($\pm\SI{0.3}{eV}$); $\chiIWO = 3.947$ / 4.228 / 4.274~eV at
2.0 / 6.3 / 13.2~nm (configuration law).

\paragraph{Source.}
\citet[Sec.~3]{Fan2021}, a composition-mixing (96/4) estimate for IWO (not bundled; recorded by the Astra
audit). Consistency check: the charge-neutrality level of \InO{} lies about 0.4~eV \emph{above} $\Ec$
\citep[pp.~503--504]{Si2021}, which is why low-barrier contacts and electron accumulation at \InO{} surfaces are
expected; the thin film's $\Ec$ rises by $\dEc$, so the neutrality level appears deeper relative to $\Ec$ in
the thinnest film.

\paragraph{Derivation (how it could be measured).}
UPS (work function and $E_F-\Ev$) combined with the optical gap, or IPES directly, gives $\chi$; an
alternative electrical route is the flat-band voltage of a MIS capacitor with a known gate metal, which
measures $\phi_M - \chi$ plus fixed charge -- exactly the degenerate combination discussed in
\cref{sec:par-qf}.

\paragraph{ATLAS implementation.}
MATERIAL AFFINITY on each IWO region \citep[p.~1337]{AtlasManual}; the confinement shift is imposed by
lowering AFFINITY by $\dEc(t)$ and raising EG300 by the same amount (\cref{sec:par-dec}).

\paragraph{Sensitivity evidence.}
\run{0020} ($\chi-0.1$~eV) and \run{0021} ($\chi+0.1$~eV) vs \run{0012}: $\Delta\Vthcc = +0.100$ / $-0.100$~V,
fixed-current SS unchanged ($+0.0\mVdec$), $\Delta\Ion = -6.9$ / $+7.0$\,\%. S6 shows that work function,
electron affinity and $\Qf$ are each an exact rigid shift at every thickness tested (max
$|\Delta\log_{10}\Id| \le 1.6\times10^{-5}$ after a 0.100~V shift), hence they cannot alter the thickness
dependence. The $\Ion$ change is the rigid shift evaluated at fixed $\Vg = 3$~V.

\paragraph{Status.}
\prov{LITERATURE\_IWO} (bulk) + \prov{DFT\_PURE\_IN2O3\_PROXY} (thickness dependence). Not identifiable from
$\Id$--$\Vg$: degenerate with $\phi_M$ and $\Qf$ (1~V $\leftrightarrow$ $5.6\times10^{12}\cmsq$, S3).

% -------------------------------------------------------------------
\section{Confinement shift of the conduction-band edge: the thickness law}
\label{sec:par-dec}

\paragraph{Definition and role.}
$\dEc(t)$ is the upward shift of the conduction-band edge of a film of thickness $t$ relative to the bulk,
attributed to geometric quantum confinement. In the classical drift--diffusion model it is inserted as a
rigid shift of the band edge (and of the gap, since $\dEc/\dEg = 1$). Because all localized-state distributions
are referenced to $\Ec$ (\cref{sec:par-tail,sec:par-deep,sec:par-dit}), the whole DOS moves with it, and the
shift acts on the transfer curve as an almost exact rigid translation in $\Vg$.

\paragraph{Values used.}
\begin{equation}
  \dEc(t) = \dEg(t) = a\,t_{\mathrm{nm}}^{-p},\qquad a = \SI{0.9205}{eV},\quad p = 1.3815,
  \label{eq:dec-law}
\end{equation}
giving $\dEc = 0.353$ / 0.072 / 0.026~eV at 2.0 / 6.3 / 13.2~nm (evidence brief).

\paragraph{Source.}
Least-squares fit of $\ln\dEc$ vs $\ln t$ through three PBE slab-minus-bulk shifts of pure \InO{}:
$0.94$~eV at 0.95~nm and $0.33$~eV at 1.98~nm (\citealp[p.~21542]{Lin2022}: slab gaps 1.88 and 1.27~eV
minus the bulk 0.94~eV), and $\dEc = 0.6$~eV for a 1.5~nm \InO{} layer on \AlO{} with $\Ev$ almost unchanged
\citep[p.~504]{Si2021}. The fit reproduces the points to $+5$ / $-12$ / $+8$\,\% (0.988, 0.526, 0.358~eV at
0.95, 1.5, 1.98~nm). The stated uncertainty is $\pm50$\,\%.

\paragraph{Derivation.} The full derivation, including the effective-mass, finite-barrier and
Schr\"odinger--Poisson cross-checks, is given in \cref{sec:der-confinement}. Two properties of the law matter
for its use. (i) It is a \emph{local} power law of three points below 2~nm; the effective-mass limit is
$t^{-2}$, so the fitted exponent 1.38 necessarily overestimates the shift when extrapolated to thicker films
($\times1.45$ at 6.3~nm and $\times2.15$ at 13.2~nm relative to the DFT-anchored effective-width estimate $E_1$ of S3,
\cref{tab:d1-e1}; $\times1.59$ / $\times2.51$ relative to the infinite-well upper bound, SYNTHESIS Q3). The two pairs
of factors describe the same overshoot against different references and do not conflict.
(ii) The law describes a subband energy, whereas a classical model needs the \emph{subthreshold-equivalent}
shift, which is smaller by the 2-D density-of-states gain (\cref{eq:d1-dec-eq}).

\paragraph{ATLAS implementation.}
MATERIAL AFFINITY$=\chi^{\mathrm{bulk}}-\dEc(t)$ and EG300$=\Eg^{\mathrm{bulk}}+\dEc(t)$ on the IWO region
\citep[pp.~1337--1338]{AtlasManual}; no Schr\"odinger or quantum-correction model is enabled in V1 (the MODELS
statement offers a Schr\"odinger--Poisson solver whose quantization direction and masses are set by SP.DIR
and the band masses, \citealp[p.~1469]{AtlasManual}; it was not used in the V1 decks).

\paragraph{Sensitivity evidence.}
Removing confinement at 2~nm (\run{0016} vs \run{0012}, same $\Qf$ and mobility) shifts $\Vthcc$ by $-0.315$~V
(0.676 $\rightarrow$ 0.361~V), $\Vthlin$ by $-0.299$~V and raises $\Ion$ by 21.4\,\%; at 13.2~nm
(\run{0017} vs \run{0013}) the shift is $-0.023$~V (evidence brief; S3). The shift is rigid: the horizontal
offset is constant at 0.345~V from $10^{-14}$ to $10^{-11}\Aum$, equal to $\dEc - \kT\ln(\Nc$ ratio$)$
$= 0.3451$~V, and the local SS vs $\log\Id$ profiles of the on/off pairs agree within 1--2$\mVdec$ (S3). The
apparent change of $\SSmin$ (73 $\rightarrow$ 84 and 151 $\rightarrow$ 139$\mVdec$) is an extraction artefact of
the floor-dependent window (\cref{sec:der-metrics}). At 6.3~nm the confinement contribution is 0.064~V
(\run{0034} minus \run{0014}, corrected for the exact $\Qf$ shift; S1).

\paragraph{Status and caveats.}
\prov{DFT\_PURE\_IN2O3\_PROXY}, extrapolated beyond 2~nm. Verdicts (S3, SYNTHESIS): film confinement at 2~nm
is physically expected (\evTHEORY; not measured on these films, and not identified from $\Id$--$\Vg$); the ATLAS value 0.353~eV (realised 0.315--0.345~V) is
\emph{plausible, unverified} and high relative to Schr\"odinger--Poisson (by 0.09~V against the central SP
case); the $t^{-1.38}$ law beyond 3~nm is \emph{rejected as physics} (kept as a numerical interpolation, its
Vth impact at 6.3~nm is $\le 0.035$~V); ``the transfer data require confinement'' is \emph{rejected}
(held-out \run{0034}, \evPOST); and $\dEc$ is \emph{not identifiable} from one $\Id$--$\Vg$ per thickness: at 2~nm
it is equivalent to $1.68$--$1.76\times10^{12}\cmsq$ of fixed negative charge (S3), i.e.\ it trades one-to-one
against $\Qf$ (\cref{sec:par-qf}; \cref{fig:symmetric_calibration}).

\begin{figure}[tbp]
  \centering
  \includegraphics[width=\linewidth]{thickness_laws}
  \caption{Thickness laws of the V1 model. (a) $\dEc(t)$: the ATLAS law $0.9205\,t^{-1.3815}$~eV
  (\cref{eq:dec-law}) against the infinite-well effective-mass $E_1$ for $\mstar = 0.18$, 0.26 and 0.35~$m_0$,
  the PBE slab points of \citet[p.~21542]{Lin2022} and \citet[p.~504]{Si2021}, and the Schr\"odinger--Poisson
  subthreshold-equivalent shifts of S3 (case Q1 at $n_s = 10^{10}\cmsq$). Above 3~nm the power law diverges
  from the $t^{-2}$ limit. (b) $\mstar(t)$ law with the PBE slab points. (c) $\Nt(t)$ power law
  against the surface\,+\,bulk alternative through the same anchors; squares: PBS devices of \citet[p.~9]{Januar2026}. (d) $\Ndeff(t)$. (e) Roughness factor
  $(1-\Dsr/t)^2$ (used in \cref{sec:par-rough}). Laws from \file{config/iwo_material_model.yaml}; DFT points
  from \file{tables/THICKNESS_LAW_EVIDENCE.csv}; SP shifts from \file{scratch/S3/s3_sp1d_notraps_out.txt}.}
  \label{fig:thickness_laws}
\end{figure}

% -------------------------------------------------------------------
\section{Derivation of the confinement shift}
\label{sec:der-confinement}

This section derives $\dEc(t)$ from first principles at increasing levels of realism and compares each level
with the law of \cref{eq:dec-law}. All numbers are from S3 (\file{S3_quantum_confinement.md} and its scripts
in \file{scratch/S3/}), which is a specialist 1-D surrogate (\evSURR) checked against ATLAS.

\paragraph{(1) Infinite square well, parabolic band.}
An electron with isotropic effective mass $\mstar$ confined between impenetrable walls at $z=0$ and $z=t$
obeys $-(\hbar^2/2\mstar)\,\psi''(z) = E\psi(z)$ with $\psi(0)=\psi(t)=0$, so $\psi_j \propto \sin(j\pi z/t)$ and
\begin{equation}
  E_j = \frac{\hbar^2\pi^2 j^2}{2\mstar t^2}
  \quad\Rightarrow\quad
  E_1 = \frac{\SI{0.376}{eV}}{(\mstar/m_0)\,t_{\mathrm{nm}}^2}.
  \label{eq:e1-infinite}
\end{equation}
For $t = 2$~nm this gives 0.522 / 0.452 / 0.366 / 0.269~eV for $\mstar = 0.18$ / 0.208 / 0.257 / 0.35~$m_0$
(S3), bracketing the ATLAS value 0.353~eV. The $t^{-2}$ scaling is exact in this limit; any fitted exponent
smaller than 2 is a sign of finite barriers and band non-parabolicity at small $t$.

\paragraph{(2) Non-parabolicity.}
In a Kane-type band $E(1+\alpha E) = \hbar^2k^2/2\mstar$ the effective mass grows with energy,
$m(E) = \mstar(1+2\alpha E)$, so a large confinement energy is partly self-limiting. Replacing $k = \pi/t$ gives
$E_1 = [-1+\sqrt{1+4\alpha E_1^{(0)}}]/(2\alpha)$ with $E_1^{(0)}$ from \cref{eq:e1-infinite}; S3 finds
0.430 / 0.380 / 0.316 / 0.240~eV at 2~nm for the same four masses. The same physics explains why the PBE slab
masses rise as the slab thins (\cref{sec:par-mstar}).

\paragraph{(3) Finite barriers.}
The film is bounded by \AlO{} (below) and air (above). With BenDaniel--Duke boundary conditions
($\psi$ and $\psi'/m$ continuous) the wave function penetrates the barriers and $E_1$ falls: S3 obtains
0.231--0.252~eV at 2~nm for a heavy barrier mass and 0.305~eV for equal masses ($\mstar = 0.18$, with
non-parabolicity), 0.039--0.040~eV at 6.3~nm and 0.010--0.011~eV at 13.2~nm.

\paragraph{(4) DFT anchoring.}
PBE slab calculations include the true band structure, the surface termination and the non-parabolicity, but
not the device environment; they underestimate absolute gaps, so only the slab-minus-bulk difference is used.
S3 constructs an ``effective width'' anchored on the slab shifts and obtains $E_1 \approx 0.35$~eV at 2~nm,
i.e.\ agreement with \cref{eq:dec-law} at 2~nm by construction (\cref{tab:d1-e1}, ratio 0.99). The slab
terminations (H-passivated/vacuum in \citealp[p.~21542]{Lin2022}; \InO{} on one \AlO{} layer with an H-terminated top surface in
\citealp[pp.~504--505]{Si2021})
differ from the device (\AlO{} below, air above), and W is absent from the slabs; hence the $\pm50$\,\%
uncertainty.

\begin{table}[tbp]
  \centering\small
  \caption{Effective-mass ground-state energy $E_1$ of the S3 effective-width model anchored on the PBE slab
  shifts, vs the ATLAS law (S3). The ratios are relative to this DFT-anchored effective-width estimate; relative to the
  infinite-well upper bound they are $\times1.59$ / $\times2.51$ at 6.3 / 13.2~nm (SYNTHESIS Q3). The law agrees at 2~nm by construction and overestimates increasingly with
  thickness because its exponent (1.38) is flatter than the effective-mass limit (2).}
  \label{tab:d1-e1}
  \begin{tabular}{lcccccc}
    \toprule
    $t$ (nm) & 2 & 3 & 4 & 5 & 6.3 & 13.2 \\
    \midrule
    $E_1$ (eV) & 0.357 & 0.187 & 0.114 & 0.076 & 0.050 & 0.012 \\
    ATLAS law, \cref{eq:dec-law} (eV) & 0.353 & 0.202 & 0.136 & 0.100 & 0.072 & 0.026 \\
    ratio law/EM & 0.99 & 1.08 & 1.19 & 1.30 & 1.45 & 2.15 \\
    \bottomrule
  \end{tabular}
\end{table}

\paragraph{(5) From a subband energy to a classical band-edge shift.}
A classical 3-D model with band edge $\Ec+\dEc$ has an electron sheet density
$n_s^{\mathrm{cl}} = \Nc t\,e^{-(\Ec+\dEc-\EF)/\kT}$ (non-degenerate), whereas the quantized film has
$n_s^{\mathrm{Q}} = \sum_i g_{2D,i}\kT\,e^{-(\Ec+E_i-\EF)/\kT}$ with $g_{2D} = \mstar/\pi\hbar^2$ per subband
(spin included). Equating the two gives the shift that a classical model must use to reproduce the
subthreshold charge,
\begin{equation}
  \dEc^{\mathrm{eq}} = -\kT\,\ln\!\left[\frac{\sum_i g_{2D,i}\,\kT\,e^{-E_i/\kT}}{\Nc\,t}\right].
  \label{eq:d1-dec-eq}
\end{equation}
Because a 2-D subband carries more states near its edge than $\Nc t$ does, $\dEc^{\mathrm{eq}} < E_1$: at 2~nm
($\mstar = 0.18$, non-parabolic, finite barrier) $\dEc^{\mathrm{eq}} = 0.207$--0.215~eV while $E_1 = 0.248$~eV;
the 2-D DOS term subtracts about 0.04~V at 2~nm and 0.01~V at 6.3~nm (S3). This is the first reason the ATLAS
rigid shift (which uses $E_1$-like DFT shifts) is on the high side.

\paragraph{(6) Self-consistent Schr\"odinger--Poisson.}
S3 solves the 1-D Schr\"odinger and Poisson equations across gate stack and film with Fermi--Dirac filling,
$\Qf = 1.73\times10^{12}\cmsq$ and the V1 $\Ndeff(t)$, trap-free (Gauss-law error $\le 2\times10^{-18}$~C/cm$^2$).
The onset shifts relative to the classical no-confinement case (C0) at $n_s = 10^{10}\cmsq$ are listed in
\cref{tab:d1-sp}; they are shown with the $n_s(\Vg)$ curves in \cref{fig:sp_confinement}. The surrogate
reproduces the ATLAS relative shifts within 11~mV (0.326~V at $10^{10}$ vs ATLAS $\Delta\Vthcc = 0.315$~V at 2~nm,
\run{0012}$-$\run{0016}; 0.023~V at 13.2~nm, \run{0013}$-$\run{0017}). The central SP case Q1 gives 0.256~V at
2~nm, the DFT-anchored width about 0.30~V and the hard wall 0.384~V; if W raises $\mstar$ to 0.26--0.35 the shift
falls to 0.14--0.20~V. The resulting bracket is 0.14--0.38~V, with 0.26--0.30~V central for $\mstar = 0.18$.

\begin{table}[tbp]
  \centering\small
  \caption{Schr\"odinger--Poisson onset shifts (V) relative to the classical no-confinement case at
  $n_s = 10^{10}\cmsq$, trap-free (S3, \file{scratch/S3/s3_sp1d_notraps.json}). C1: classical with the ATLAS law;
  Q1: central quantum case ($\mstar = 0.18$, finite barriers, non-parabolic); Q2: hard wall; Q3: BenDaniel--Duke
  with heavy barrier mass; Q4/Q5: $\mstar = 0.257$ / 0.35.}
  \label{tab:d1-sp}
  \begin{tabular}{lcccccc}
    \toprule
    $t$ (nm) & C1 & Q1 & Q2 & Q3 & Q4 & Q5 \\
    \midrule
    2.0  & 0.345 & 0.256 & 0.384 & 0.202 & 0.205 & 0.139 \\
    3.0  & 0.197 & 0.131 & 0.179 & 0.105 & 0.090 & 0.052 \\
    4.0  & 0.133 & 0.078 & 0.102 & 0.065 & 0.047 & 0.021 \\
    5.0  & 0.097 & 0.052 & 0.065 & 0.045 & 0.028 & 0.007 \\
    6.3  & 0.070 & 0.036 & 0.043 & 0.031 & 0.016 & $-0.002$ \\
    13.2 & 0.026 & 0.016 & 0.018 & 0.015 & 0.001 & $-0.012$ \\
    \bottomrule
  \end{tabular}
\end{table}

\begin{figure}[tbp]
  \centering
  \includegraphics[width=\linewidth]{sp_confinement}
  \caption{S3 1-D Schr\"odinger--Poisson surrogate (trap-free, \file{scratch/S3/s3_sp1d_notraps.json}): electron
  sheet density $n_s(\Vg)$, quantum vs classical, at 2~nm and 6.3~nm, with the onset shift and the electron
  centroid. Shifts vs the classical case C0 at $10^{10}\cmsq$: C1 (ATLAS law) 0.345 / 0.070 / 0.026~V and
  Q1 0.256 / 0.036 / 0.016~V at 2 / 6.3 / 13.2~nm. The shift is rigid at 2~nm, and at 6.3~nm the
  law exceeds the quantum shift by about a factor 2.}
  \label{fig:sp_confinement}
\end{figure}

\paragraph{(7) Field quantization and capacitance.}
Above threshold the gate field $F$ forms a triangular well at the front interface, whose ground state in the
triangular-well approximation is $E_1^{\mathrm{tri}} \simeq (\hbar^2/2\mstar)^{1/3}(9\pi qF/8)^{2/3}$. S3 finds
$E_1^{\mathrm{tri}} = 0.10$~eV and centroid 3.45~nm at $n_s = 10^{12}\cmsq$ ($F = 1.95\times10^5$~V/cm), and
0.56~eV with centroid 1.47~nm at $n_s = 1.3\times10^{13}\cmsq$ (the bound $\Cox(3-\Vthcc)/q$ at $\Vg = 3$~V). The
2~nm film is never field-quantized (centroid 0.94--0.98~nm $\approx t/2$ up to $1.3\times10^{13}\cmsq$); the 6.3 and
13.2~nm films are, above about $2$--$3\times10^{12}$ and $1$--$3\times10^{11}\cmsq$ respectively. The effective
capacitance $C_{\mathrm{eff}}/\Cox$ over 2--3~V is already 0.87--0.88 classically (centroid and degenerate DOS),
and quantum dark space lowers it by $<1$\,\% at 2~nm and 3--5\,\% at 6.3 and 13.2~nm (S3). This on-state
field quantization is not represented in V1; its implication for the mobility is discussed in
\cref{sec:par-mu}.

\paragraph{(8) Two neglected corrections.}
\emph{Image (dielectric) confinement}: the air side ($\varepsilon = 1$) against IWO ($\varepsilon = 9.3$) adds
about +31~meV at 1~nm from the surface, roughly +0.02--0.03~eV at 2~nm (S3 estimate; \emph{plausible,
unverified}). \emph{Energy reference of the tail}: if the localized tail states follow the confined edge the
2~nm onset shift is 0.26~V (Q1 with traps); if they stay at the bulk edge (localization length shorter than $t$)
the whole $4\times10^{12}\cmsq$ tail must fill before the subband populates and the onset moves by 0.91~V (S3).
This is the dominant model-form uncertainty of the confinement treatment and is \emph{not determined} from the
available data. With local tails, field quantization at 6.3 and 13.2~nm raises $\Vg(n_s = 10^{12})$ by
0.15--0.19~V relative to classical while $\Vg(10^{10})$ moves by $\le 0.04$~V -- a change of curve shape the rigid
law cannot represent (S3; \evSURR, registered as a prediction).

\paragraph{Summary of the derivation.}
The law of \cref{eq:dec-law} is consistent with effective-mass theory at 2~nm (inside the 0.14--0.38~V bracket)
but on its high side, and it overestimates the S3 subthreshold-equivalent shift by about $2\times$ at 6.3~nm (0.070 vs 0.036~V) and $1.6\times$ at
13.2~nm (0.026 vs 0.016~V). The shift could be measured through the thickness dependence of the optical gap
(expected +0.2--0.45~eV at 2~nm vs 13.2~nm), through IPES or UPS/XPS for the CBM, and through a denser thickness series
(2$\rightarrow$3~nm: $\Delta\Vthcc$ +0.125--0.205~V predicted; S3 falsification table).

% -------------------------------------------------------------------
\section{Conduction-band effective mass}
\label{sec:par-mstar}

\paragraph{Definition and role.}
$\mstar(t)$ is the density-of-states effective mass of the conduction band. In the V1 decks (no
Schr\"odinger solver, constant mobility) it enters \emph{only} through the effective density of states
$\Nc \propto \mstar{}^{3/2}$ (\cref{sec:par-nc}); its influence on the mobility (Drude $\mu = q\tau/\mstar$) is absorbed
into the fitted band mobility (\cref{sec:par-mu}), and its influence on the confinement energy is already
contained in the DFT-based $\dEc$ law.

\paragraph{Values used.}
\begin{equation}
  \mstar(t) = m_{\mathrm{bulk}} + b\,t_{\mathrm{nm}}^{-q},\qquad m_{\mathrm{bulk}} = 0.208\,m_0,\quad
  b = 0.1311\,m_0,\quad q = 1.4121,
  \label{eq:d1-mstar}
\end{equation}
giving $\mstar = 0.257$ / 0.218 / 0.211~$m_0$ at 2.0 / 6.3 / 13.2~nm (evidence brief).

\paragraph{Source.}
Bulk value: optical-Hall measurement on a single \InO{} crystal ($n = 2.8\times10^{17}\cmcu$),
$0.208\,m_0$ \citep[p.~225102-1, abstract]{Stokey2021}; the zero-density extrapolation quoted there is $0.18\,m_0$
\citep[p.~225102-2]{Stokey2021}. Thickness increment: PBE slab masses of pure \InO{}, 0.17 (bulk), 0.19
(3.52~nm), 0.23 (1.98~nm) and 0.30~$m_0$ (0.95~nm) \citep[p.~21540]{Lin2022}. The effect of W on $\mstar$ is
reported only qualitatively in the target paper (flatter CBM, heavier mass; text lines 238--241 as quoted in S3);
it is \emph{not determined} numerically and is not included.

\paragraph{Derivation.}
Only PBE \emph{differences} are used, as for $\dEc$: $\Delta m(t) = m_{\mathrm{slab}}(t) - m_{\mathrm{PBE,bulk}}$
$= 0.13$, 0.06 and 0.02~$m_0$ at 0.95, 1.98 and 3.52~nm. A least-squares fit of $\ln\Delta m$ vs $\ln t$ gives
$b$ and $q$ of \cref{eq:d1-mstar}; the fitted increments are 0.141 / 0.050 / 0.022~$m_0$ at the three slab
thicknesses. The increment is added to the \emph{experimental} bulk mass. Physically the increase is the
non-parabolicity of \cref{sec:der-confinement}: a confined electron sits at a higher kinetic energy, where
$m(E) = \mstar(1+2\alpha E)$ is heavier. The mass could be measured by optical Hall or infrared
ellipsometry (plasma frequency with known $n$) on films of each thickness, or from the temperature dependence of
Shubnikov--de Haas oscillations in thick films.

\paragraph{ATLAS implementation.}
Not a direct deck parameter in V1: $\mstar$ is converted into NC300 of the MATERIAL statement
\citep[p.~1340]{AtlasManual} (\cref{sec:par-nc}). The band masses of the MATERIAL statement would matter only if
the Schr\"odinger model (SP.DIR masses, \citealp[p.~1469]{AtlasManual}) were enabled, which it is not.

\paragraph{Sensitivity evidence.}
$\mstar\times1.3$ (isolation tests, DOS 384/192, vs the same-mobility baselines): $\Delta\Vthcc = -0.044$~V
(2~nm, \run{0050} vs \run{0038}) and $-0.029$~V (13.2~nm, \run{0051} vs \run{0040}); $\SScc$ $-17.7$ / $-13.2\mVdec$;
$\Delta\Ion = +4.1$ / $+5.8$\,\% (S6; SYNTHESIS). The shift is \emph{not} rigid (0.17~dec change of shape at
2~nm, S6), in contrast to the work-function and affinity shifts. The evidence brief quotes $-0.028$~V and
$+3.7$\,\% for \run{0051}; S6 and SYNTHESIS, which are later and state the baseline explicitly, are used here. In
the budget, $\mstar(t)$ contributes $+14.9$~mV to the 2$\rightarrow$13.2~nm $\Vthcc$ step (S6). The reason a
30\,\% heavier mass (only $\kT\ln 1.3^{3/2} = \SI{10}{mV}$ of Fermi-level shift at fixed free density) moves
$\Vthcc$ by 44~mV is the tail: at fixed free density a lower $\EF$ empties part of the tail charge, and the
tail term is ten times larger than the free-carrier term at threshold (\cref{tab:d1-budget}).

\paragraph{Status.}
\prov{MEASURED\_RELATED\_MATERIAL} (bulk, pure \InO{}) + \prov{DFT\_PURE\_IN2O3\_PROXY} (increment).
Non-parabolicity is otherwise ignored; W-induced mass change \emph{not determined}.

% -------------------------------------------------------------------
\section{Effective density of states of the conduction band}
\label{sec:par-nc}

\paragraph{Definition and role.}
$\Nc$ sets the free-electron density for a given $\EF - \Ec$,
$n = \Nc\,\mathcal{F}_{1/2}[(\EF-\Ec)/\kT]$ (Fermi--Dirac statistics, $\mathcal{F}_{1/2}$ the normalised
Fermi integral). It fixes where the Fermi level sits at threshold and therefore the partition between free
and trapped electrons.

\paragraph{Values used.}
\begin{equation}
  \Nc = 2\left(\frac{2\pi\mstar\kT}{h^2}\right)^{3/2}
      = \SI{2.51e19}{cm^{-3}}\left(\frac{\mstar}{m_0}\right)^{3/2}\left(\frac{T}{\SI{300}{K}}\right)^{3/2},
  \label{eq:d1-nc}
\end{equation}
which with \cref{eq:d1-mstar} evaluates to $3.28$ / $2.55$ / $2.44\times10^{18}\cmcu$ at 2.0 / 6.3 / 13.2~nm
($2.38\times10^{18}\cmcu$ for the bulk mass 0.208).

\paragraph{Source and derivation.}
\Cref{eq:d1-nc} is the standard result for a parabolic 3-D band: the number of states per volume below
energy $E$ above the edge is $(1/3\pi^2)(2\mstar E/\hbar^2)^{3/2}$ per spin, so $g(E) = (1/2\pi^2)(2\mstar/\hbar^2)^{3/2}\sqrt{E}$
including spin; integrating $g(E)e^{-E/\kT}$ gives $\Nc$. At 2~nm the film is quasi-2D and $\Nc$ is an
effective continuum parameter; the correct 2-D occupation is the subject of \cref{eq:d1-dec-eq}. The ratio
$\Nc(2\,\mathrm{nm})/\Nc(\mathrm{bulk}) = 1.376$ corresponds to $\kT\ln 1.376 = \SI{8}{mV}$ of Fermi-level shift, the
value S1 quotes for the analytic $\kT\ln\Nc$ term at 2~nm. The same ratio explains why removing confinement
(which also removes the mass increment) shifts the 2~nm curve by $\dEc - \kT\ln(\Nc$ ratio$) = 0.345$~V rather
than 0.353~V (S3).

\paragraph{ATLAS implementation.}
MATERIAL NC300 per IWO region \citep[p.~1340]{AtlasManual}; Fermi--Dirac statistics are enabled on the MODELS
statement (FERMIDIRAC, \citealp[p.~1433]{AtlasManual}). NC300 is the 300~K value; its temperature scaling in the
elevated-temperature predictions is treated in \cref{sec:par-temperature}.

\paragraph{Sensitivity evidence.}
Through $\mstar$: runs \run{0050}/\run{0051} (\cref{sec:par-mstar}). In the 1-D surrogate, ``$\mstar$ through $\Nc$
only'' contributes $-0.039$ / $-0.008$ / $-0.002$~V to $\Vthcc$ at 2 / 6.3 / 13.2~nm (S1). REVIEW (O2) notes that the
$\Nc$ part of the confinement removal gives 43~mV and $14.6\mVdec$, so the confinement degeneracy with $\Qf$ is
near-exact, not exact.

\paragraph{Status.}
\prov{CALCULATED} from $\mstar$ (itself a proxy). Caveat: continuum 3-D DOS used for a quasi-2-D film.

% -------------------------------------------------------------------
\section{Effective density of states of the valence band}
\label{sec:par-nv}

\paragraph{Definition and role.}
$\Nv$ enters only the intrinsic density $n_i = \sqrt{\Nc\Nv}\,e^{-\Eg/2\kT}$, because the model solves
electrons only (holes are not transported).

\paragraph{Values used.}
$\Nv = 1.0\times10^{19}\cmcu$ at 300~K, all thicknesses; hole mobility $0.01\cmVs$ (both \prov{ASSUMED}).

\paragraph{Source and derivation (how it could be derived).}
No source; the value is a generic placeholder. With the values of \cref{sec:par-eg,sec:par-nc},
$n_i \approx 2\times10^{-10}$ (2~nm) to $7\times10^{-8}\cmcu$ (13.2~nm) (evaluated from the formula), so no
physically plausible $\Nv$ ($10^{18}$--$10^{20}\cmcu$) can influence an $n$-type transfer curve. If needed,
$\Nv$ would follow from \cref{eq:d1-nc} with the valence-band DOS mass, which for \InO{} is large and poorly
known; DFT hole masses of the slabs could supply it.

\paragraph{ATLAS implementation.}
MATERIAL NV300 \citep[p.~1340]{AtlasManual}.

\paragraph{Sensitivity evidence.}
None run; insensitive by construction (\evTHEORY).

\paragraph{Status.}
\prov{ASSUMED}; irrelevant to the fitted curves.

% -------------------------------------------------------------------
\section{Acceptor-like band-tail states}
\label{sec:par-tail}

\paragraph{Definition and role.}
The exponential tail of localized, acceptor-like states below the conduction-band edge is the dominant
charge reservoir of the model near threshold. When filled, the states are negatively charged; they pin the
Fermi level, set the subthreshold slope in the $10^{-10}$--$10^{-8}\Aum$ window and the trapped/free partition
above threshold. The density of states is
\begin{equation}
  g_{\mathrm{TA}}(E) = \NTA\,\exp\!\left(\frac{E-\Ec}{\WTA}\right),\quad E<\Ec,
  \qquad
  \int_{\Ec-\Eg}^{\Ec} g_{\mathrm{TA}}\,dE = \NTA\WTA\left(1-e^{-\Eg/\WTA}\right) \simeq \NTA\WTA \equiv \Nt .
  \label{eq:tail-dos}
\end{equation}
The model is parameterised by the integral $\Nt(t)$ and the width $\WTA$, and $\NTA = \Nt/\WTA$.

\paragraph{Values used.}
$\Nt(t) = \Nt^{(2)}(2/t_{\mathrm{nm}})^{\alpha}$ with $\Nt^{(2)} = 2\times10^{19}\cmcu$ and $\alpha = 0.75$;
$\WTA = \SI{0.040}{eV}$ shared (tail temperature $T_t = \WTA/\kB = \SI{464}{K}$). Per film:
\begin{center}\small
\begin{tabular}{lccc}
  \toprule
  $t$ (nm) & 2.0 & 6.3 & 13.2 \\
  \midrule
  $\Nt$ ($\cmcu$) & $2.0\times10^{19}$ & $8.5\times10^{18}$ & $4.9\times10^{18}$ \\
  $\NTA = \Nt/\WTA$ ($\si{cm^{-3}.eV^{-1}}$) & $5.0\times10^{20}$ & $2.1\times10^{20}$ & $1.2\times10^{20}$ \\
  sheet $\Nt t$ ($\cmsq$) & $4.0\times10^{12}$ & $5.3\times10^{12}$ & $6.4\times10^{12}$ \\
  filled at $\Vthcc$ ($\cmsq$; S1, 1-D) & $1.43\times10^{12}$ & $1.27\times10^{12}$ & $5.5\times10^{11}$ \\
  $q n_t/\Cox$ at $\Vthcc$ (V; S1) & 0.256 & 0.226 & 0.098 \\
  \bottomrule
\end{tabular}
\end{center}
(6.3~nm columns for the validation configuration \run{0014}.)

\paragraph{Source.}
$\Nt^{(2)}$ is \prov{FITTED}: the product $\NTA\WTA$ of the 2~nm calibration in ATLAS (\cref{ch:calibration}).
The exponent comes from a two-point ratio of about 4 between the 13.2 and 2~nm trap densities of the target
paper's thickness series \citep[Sec.~2.8, Fig.~5a, p.~7]{Januar2026}: $\alpha = \ln 4/\ln(13.2/2) = 0.73 \approx 0.75$.
The same paper's compact-model values for its bias-stress devices, $5.3\times10^{19}$ (2~nm) and
$3.39\times10^{18}\cmcu$ (10~nm) \citep[p.~9]{Januar2026}, use a different ($\theta_t$) convention and different
devices; they imply a ratio 15.6 and a sheet density that \emph{decreases} with $t$, which forces a negative
bulk term and is unphysical for this series (S4), so they are not used as anchors. Context: the conductance-method
bulk DOS of ALD \InO{} is $3.3\times10^{20}\si{cm^{-3}.eV^{-1}}$ at $\Nd \approx 10^{20}\cmcu$
\citep[pp.~4--5]{Wang2022}. $\WTA$ is fitted on the 2~nm subthreshold slope; the target paper reports only the
stress-induced change of the tail temperature, so the absolute $T_t$ is \emph{not determined} from the literature.

\paragraph{Derivation.}
\emph{Occupation.} The trapped density at quasi-Fermi level $u = E_{Fn}-\Ec$ is
$n_t(u) = \int g_{\mathrm{TA}}(E)\,f(E-u)\,dE$ with the Fermi function $f$ (\cref{sec:par-capture} shows why the
steady-state occupancy is Fermi--Dirac regardless of capture cross sections). For $u$ several $\kT$ below $\Ec$
and $\WTA > \kT$, extending the integral to $+\infty$ gives
\begin{equation}
  n_t(u) \simeq \NTA\WTA\,e^{u/\WTA}\;\frac{\pi\kT/\WTA}{\sin(\pi\kT/\WTA)},
  \label{eq:d1-nt-occ}
\end{equation}
where the last factor is the thermal enhancement over a step-function occupation; for $\WTA = 40$~meV it is
2.26 (the value used by S4). Near $\Ec$ the truncation of the tail at $\Ec$ reduces it. \Cref{eq:d1-nt-occ} shows
the two regimes of the subthreshold: when $\WTA > \kT$ the trapped charge grows more slowly with $u$ than the
free charge ($e^{u/\WTA}$ vs $e^{u/\kT}$), so it dominates at low current and yields to the free charge at high
current. Its trap capacitance $C_t = q^2 t\,\partial n_t/\partial u = q^2 t\,n_t/\WTA$ enters the slope
(\cref{eq:ss-dit}).
\emph{Thickness law.} The power law is an interpolation. S4 shows that the surface\,+\,bulk form
$\Nt t = N_s + N_b t$ through the same anchors ($N_s = 3.57\times10^{12}\cmsq$, $N_b = 2.15\times10^{18}\cmcu$)
predicts $\Nt(6.3) = 7.8\times10^{18}$, 8\,\% below the power law, and diverges only beyond the fitted range
(31.8~nm: $3.3$ vs $2.5\times10^{18}\cmcu$). The trap-DOS inversion (\cref{sec:der-ss}) supports the surface\,+\,bulk
form: $D(\Ec-0.1\,\mathrm{eV}) = 1.05$ / 1.67 / $2.16\times10^{13}\si{cm^{-2}.eV^{-1}}$, i.e.\ $\propto t^{0.38\pm0.15}$,
consistent with $D_s \approx 8.5\times10^{12}\si{cm^{-2}.eV^{-1}}$ plus $g_b \approx 10^{19}\si{cm^{-3}.eV^{-1}}$ (S4).
Suitable measurements are temperature-dependent transfer curves (activation energy vs $u$), quasi-static
C--V (free/trapped partition), photo-excited charge collection, and Hall measurements on unpatterned films.

\paragraph{ATLAS implementation.}
DEFECTS CONTINUOUS with NTA (the density at the conduction-band edge, \citealp[p.~1171]{AtlasManual}) and WTA
(the characteristic decay energy of the acceptor tail, \citealp[p.~1173]{AtlasManual}); the DEFECTS statement
admits one exponential and one Gaussian distribution per state type \citep[p.~1169]{AtlasManual}. The continuous
distribution is discretised on NUMA/NUMD energy levels (default 12, \citealp[p.~1168]{AtlasManual}); the V1 decks
use 96/48 and, from campaign~B, 384/192 (\cref{sec:der-dos-discretization}). The tail capture cross section is
$10^{-15}\si{cm^2}$ (\cref{sec:par-capture}).

\paragraph{Sensitivity evidence.}
$\Nt\times1.5$ (\run{0023} vs \run{0012}): $\Delta\Vthcc = +0.082$~V, fixed-current SS ($10^{-10}$--$10^{-9}\Aum$)
$+49.0\mVdec$, $\Delta\Ion = -25.1$\,\%. $\WTA + 5$~meV (\run{0024}): $+0.020$~V, $+4.6\mVdec$ in the same window
($\SScc$ $+12\mVdec$, evidence brief), $\Delta\Ion = +0.4$\,\%. The tail is the largest non-confinement term of the
threshold budget (0.10--0.26~V, 25\,\% of the 2$\rightarrow$13.2~nm step; S1). The pre-registered S2
re-partition with $\Nt(6.3) = 2.66\times10^{19}\cmcu$ (B0, \run{0037}) met the on-state targets but gave
$\SScc = 410.9$ vs $295.4\mVdec$ measured and was \emph{rejected} (\evOOS, failed; SYNTHESIS). A broader tail
($\WTA$ 60~meV) was disfavoured in the S4 surrogate ($\SScc$ +126$\mVdec$).

\paragraph{Status.}
\prov{FITTED} ($\Nt^{(2)}$, $\WTA$) with a \prov{LITERATURE\_IWO} trend. Verdicts (S4): an exponential acceptor tail
near $\Ec$ controls threshold and $\SScc$ in all films -- \emph{strongly supported} (inversion ratio
measured/ATLAS 0.81--1.17 for $u = -0.125$ to $-0.025$~eV at 2 and 13.2~nm); $\Nt\propto t^{-0.75}$ as physics --
\emph{rejected} (kept as interpolation); surface\,+\,bulk -- \emph{plausible, unverified} (preferred form).
Caveat: $\Nt$, $\WTA$ and $\Qf$ trade off on a single curve.

% -------------------------------------------------------------------
\section{Deep acceptor Gaussian}
\label{sec:par-deep}

\paragraph{Definition and role.}
A Gaussian band of acceptor-like states deep in the gap,
\begin{equation}
  g_{\mathrm{GA}}(E) = \NGA\exp\!\left[-\left(\frac{E - (\Ec-\EGA)}{\WGA}\right)^{2}\right],\qquad
  \int g_{\mathrm{GA}}\,dE = \sqrt{\pi}\,\NGA\WGA,
  \label{eq:d1-gauss}
\end{equation}
representing deep (e.g.\ oxygen-vacancy related) states. Below threshold these states are full, so they act as
a fixed negative charge; they matter only through their total number.

\paragraph{Values used.}
$\NGA = 5\times10^{16}\si{cm^{-3}.eV^{-1}}$, $\EGA = 0.6$~eV below $\Ec$, $\WGA = 0.15$~eV, all thicknesses. Total
density $\sqrt{\pi}\NGA\WGA = 1.33\times10^{16}\cmcu$, i.e.\ sheets of $2.7\times10^{9}$ / $8.4\times10^{9}$ /
$1.75\times10^{10}\cmsq$ at 2.0 / 6.3 / 13.2~nm and threshold contributions $\le 0.003$~V (\cref{eq:d1-qcox}),
matching S1's 0.000 / 0.001 / 0.003~V.

\paragraph{Source.}
\prov{ASSUMED}; the magnitude is within the order of the range $2.5$--$3.7\times10^{16}$ reported by \citet{Fan2021}
for a different IWO device (not bundled; value recorded in the configuration). Donor-like tail and donor Gaussian
are set to zero (valence tail negligible for $n$-type transfer curves; the donor background is represented by
$\Ndeff$, \cref{sec:par-nd}).

\paragraph{Derivation (how it could be measured).}
Photo-excited C--V or sub-gap photocurrent spectroscopy resolve deep states; the transfer curve is blind to them
because they never change occupancy in the swept range. The trap-DOS inversion of \cref{sec:der-ss} does see the
states at $\Ec-0.15$ to $-0.2$~eV and finds that V1 underestimates them by $\times2.2$--2.8 (2~nm) and $\times2.6$
(6.3~nm) (S4), consistent with the model slope being too soft in the $10^{-10}$--$10^{-9}\Aum$ window
(\cref{sec:cal-final}); the lower $\SSmin$ of \run{0012} (73 vs 84.5$\mVdec$) is not used as evidence, because
$\SSmin$ is not reliable (\cref{sec:der-metrics}).

\paragraph{ATLAS implementation.}
DEFECTS NGA, EGA, WGA (acceptor Gaussian; NGA is the density of acceptor-like states in the Gaussian,
\citealp[p.~1171]{AtlasManual}; defaults NGA $5\times10^{17}\si{cm^{-3}.eV^{-1}}$, WGA 0.1~eV,
\citealp[pp.~1168--1169]{AtlasManual}).

\paragraph{Sensitivity evidence.}
No run varies the deep Gaussian (insensitive by the estimate above). A \emph{near-edge} acceptor Gaussian was tested
as a hypothesis for the 6.3~nm on-state shape (campaign C, \run{0052}): $2\times10^{12}\cmsq$ at $\Ec-0.03$~eV,
$\WGA = 0.02$~eV (i.e.\ $\NGA = 2\times10^{12}/(6.3\times10^{-7}\cdot0.02\sqrt{\pi}) = 9.0\times10^{19}\si{cm^{-3}.eV^{-1}}$
by \cref{eq:d1-gauss}), with $\muband = 15.4\cmVs$. Against the S4 prediction (gap $+0.28$~V, $\gm$ $+24$\,\%, $\SScc$
$+23\mVdec$, $\Ion \pm3$\,\%) it gave gap $+0.119$~V, $\gm$ $+25.5$\,\%, $\SScc$ $+15.6\mVdec$ and $\Ion$ $+16.0$\,\%:
\emph{partly failed} (\evOOS; S6, SYNTHESIS), recovering only 33\,\% of the gap (\cref{ch:6p3}).

\paragraph{Status.}
\prov{ASSUMED}; insensitive (\evTHEORY). Deeper states are underestimated (\emph{strongly supported}, S4).

% -------------------------------------------------------------------
\section{Front interface acceptor sheet and its regularisation}
\label{sec:par-dit}

\paragraph{Definition and role.}
Acceptor-like interface states at the IWO/\AlO{} interface, described as a Gaussian in energy per unit area,
$\Dit(E) = D_{\mathrm{it}}^{\mathrm{pk}}\exp\{-[(E-(\Ec-E_0))/W_{\mathrm{it}}]^2\}$. They add a sheet charge that is
full below threshold and a capacitance $C_{\mathrm{it}} = q^2\Dit$ where they cross the Fermi level.

\paragraph{Values used.}
$D_{\mathrm{it}}^{\mathrm{pk}} = 3\times10^{11}\si{cm^{-2}.eV^{-1}}$, $E_0 = 0.3$~eV, $W_{\mathrm{it}} = 0.12$~eV, shared across
thickness. Total sheet $\sqrt{\pi}D_{\mathrm{it}}^{\mathrm{pk}}W_{\mathrm{it}} = 6.4\times10^{10}\cmsq$.

\paragraph{Source.}
\prov{ASSUMED}. Literature for the \HfO/\InO{} interface: $6\times10^{11}\si{cm^{-2}.eV^{-1}}$ from an SS of
63.5$\mVdec$ \citep[p.~21539]{Lin2022} and $6.3\times10^{11}$ from the subthreshold method \citep[pp.~4--5]{Wang2022}.
The device interface here is IWO/\AlO, for which no value is available; the assumed peak is half the \HfO/\InO{}
proxy.

\paragraph{Derivation: regularisation of a sheet.}
ATLAS defect distributions are volumetric (per cm$^3$). A sheet density is represented by a thin IWO layer of
thickness $d_{\mathrm{reg}}$ adjacent to the interface carrying the volume density
\begin{equation}
  N_{\mathrm{GA}}^{\mathrm{vol}} = \frac{D_{\mathrm{it}}^{\mathrm{pk}}}{d_{\mathrm{reg}}},\qquad
  d_{\mathrm{reg}} = \SI{0.25}{nm}\;\Rightarrow\;
  N_{\mathrm{GA}}^{\mathrm{vol}} = \frac{3\times10^{11}}{2.5\times10^{-8}} = 1.2\times10^{19}\,\si{cm^{-3}.eV^{-1}}.
  \label{eq:dit-regularization}
\end{equation}
The representation is exact in the limit $d_{\mathrm{reg}}\rightarrow0$, and its error is of order the potential drop
across the layer, $\Delta\psi \approx q N_{\mathrm{it}} d_{\mathrm{reg}}/2\varepsilon_s$, which must be small compared
with $\kT/q$: with the full sheet $6.4\times10^{10}\cmsq$, $\Delta\psi \approx 1.6\times10^{-19}\cdot6.4\times10^{10}\cdot2.5\times10^{-8}/(2\cdot8.2\times10^{-13})
\approx 0.2$~mV. The condition is equivalently $d_{\mathrm{reg}} \ll L_D$ (\cref{sec:der-depletion}). Numerical check:
halving the layer changes $\Id$ by $\le 0.009$~dec and $-0.06$\,\% (S6 check E; configuration note). At threshold
($\EF$ about 0.04~eV below $\Ec$, S1) the whole Gaussian at $\Ec-0.3$~eV is full, giving
$q\cdot6.4\times10^{10}/\Cox = 0.011$~V, consistent with S1's 0.012~V.
The interface-state density could be measured by conductance or charge-pumping methods on the target stack, or by the trap-DOS
inversion of \cref{sec:der-ss} at energies where the tail is small; the SS formula (\cref{eq:ss-dit}) gives only an
upper bound when bulk traps are present.

\paragraph{ATLAS implementation.}
A 0.25~nm IWO region at the \AlO{} interface with a DEFECTS acceptor Gaussian (NGA, EGA, WGA;
\citealp[p.~1171]{AtlasManual}) of amplitude \cref{eq:dit-regularization}.

\paragraph{Sensitivity evidence.}
$\Dit\times3$ (\run{0025} vs \run{0012}): $\Delta\Vthcc = +0.025$~V, fixed-current SS $+0.1\mVdec$ ($10^{-10}$--$10^{-9}$;
$\SScc$ $+9\mVdec$, evidence brief), $\Delta\Ion = -1.8$\,\%. $\Dit\times5$ at 6.3~nm (A4, \run{0033} vs \run{0014}):
$+0.051$~V, gap $\Vthlin-\Vthcc$ unchanged (0.857 vs 0.862~V; the $\SSmin$ change, $+8.3\mVdec$, is quoted for
history only). About 0.012--0.013~V per multiple of the assumed value; closing the 0.29~V 6.3~nm miss would need a
peak of about $7\times10^{12}$ ($\times23$--24) at 6.3~nm only, on a gate stack shared with the other films, so front
$\Dit$ is \emph{rejected} as that cause (S1, S4; SYNTHESIS~D; \evPOST).

\paragraph{Status.}
\prov{ASSUMED} (value), \prov{CALCULATED} (regularisation, \evNUM-verified). Interface-trap charge
\emph{demonstrated small} (0.012~V, thickness-independent).

% -------------------------------------------------------------------
\section{Fixed interface charge}
\label{sec:par-qf}

\paragraph{Definition and role.}
$\Qf$ is a fixed positive sheet charge (per cm$^2$) at the front IWO/\AlO{} interface. It is the model's
alignment parameter: it shifts the whole transfer curve rigidly along $\Vg$ and absorbs every other rigid offset
(gate work function, affinity, $\Cox$ uncertainty, fixed oxide charge, ionised donors within about 1~nm of the
interface).

\paragraph{Values used.}
$\Qf = 1.73\times10^{12}\cmsq$, shared by 2.0 and 13.2~nm (and by the 6.3~nm validation configuration
\run{0014}); $8.7\times10^{10}\cmsq$ for the device-specific ``tuned'' 6.3~nm configuration (\run{0015}, and the
DOS~384/192 recalibration \run{0039}) (evidence brief; predictions register).

\paragraph{Source.}
\prov{FITTED}: stage~3 of the calibration (\cref{sec:cal-history}), one value shared across thickness. No
literature value exists for this interface.

\paragraph{Derivation.}
Gauss's law across the gate dielectric, with the sheet at the semiconductor side of the interface, requires the
gate to image $-q\Qf$; for fixed channel potential this is a flat-band shift
\begin{equation}
  \Delta V_{\mathrm{FB}} = -\frac{q\,\Qf}{\Cox},
  \label{eq:flatband-shift}
\end{equation}
i.e.\ $-0.310$~V for $1.73\times10^{12}$ and $-0.016$~V for $8.7\times10^{10}\cmsq$ (\cref{eq:d1-qcox}). Because the
charge sits at the interface it does not alter the potential profile inside the film, so the shift is exact at every
$\Vg$ (confirmed to 1~mV in ATLAS, \cref{sec:par-stack}). The charge is modelled as a sheet rather than a volume for the following reason. Spread through the film,
$1.73\times10^{12}\cmsq$ corresponds to $8.7\times10^{18}$ / $2.7\times10^{18}$ / $1.3\times10^{18}\cmcu$; a
\emph{volumetric} donor of $8.7\times10^{18}$ would add about $-2.0$~V at 13.2~nm (\cref{eq:d1-nd-shift}), so the shared
fit requires the charge to be sheet-like at the front. It is electrostatically identical to a layer of ionised oxygen
vacancies of about $1.7\times10^{19}\cmcu$ within 1~nm of the \AlO{} (S4) -- the two are indistinguishable on
$\Id$--$\Vg$. The required $\Qf$ is degenerate with, and therefore depends on, the confinement assumption. Per film, with confinement
the values that would fit each device alone are $1.81$ / $0.09$ / $1.56\times10^{12}\cmsq$ and without confinement
$0.047$ / $-0.27$ / $1.43\times10^{12}\cmsq$ (S1; SYNTHESIS D3; \cref{fig:symmetric_calibration}); the best single $\Qf$ over the
three films is $1.153\times10^{12}$ (rms residual 0.136~V) with and $0.403\times10^{12}$ (0.133~V) without confinement
(SYNTHESIS D3). Both readings need one film-specific charge of about $1.4$--$1.6\times10^{12}\cmsq$. The charge could
be measured through the flat-band voltage of MIS capacitors with several EOTs of the same stack, $V_{\mathrm{FB}} =
(\phi_M-\chi-\ldots)/q - q\Qf\,\EOT/(3.9\varepsilon_0)$, separates $\Qf$ (slope) from the work-function difference
(intercept).

\paragraph{ATLAS implementation.}
INTERFACE QF$=1.73\times10^{12}$ at the IWO/\AlO{} interface; QF is given in electronic charges per cm$^2$, a positive
value introduces positive charge \citep[p.~1260]{AtlasManual}; the statement applies by default to
semiconductor--insulator interfaces \citep[p.~1264]{AtlasManual} (example of use, \citealp[p.~1263]{AtlasManual}).

\paragraph{Sensitivity evidence.}
$\Qf$ $0\rightarrow1.73\times10^{12}$ at 2~nm (\run{0001}$\rightarrow$\run{0003}): $\Delta\Vthcc = -0.309$~V,
fixed-current SS $+0.4\mVdec$, $\Delta\Ion = +27.1$\,\%; at 13.2~nm \run{0002}$\rightarrow$\run{0004}. S6: an exact
rigid shift (max $|\Delta\log_{10}\Id| \le 1.6\times10^{-5}$ for the equivalent 0.1~V work-function shift). At 6.3~nm
the tuned $\Qf$ (\run{0015}) fixes $\Vthcc$ ($+0.007$~V) and $\Ion$ but leaves $\Vthlin$ $-0.35$~V and $\gmmax$ $-23$\,\%
(\cref{ch:6p3}).

\paragraph{Status.}
\prov{FITTED}, classified as an effective electrostatic correction. REVIEW: with $\phi_M$, $\chi$, $\Dit$, the deep
Gaussian and $\varepsilon_{\mathrm{Al_2O_3}}$ all inside degenerate sets, $\Qf$ is ``an alignment absorber, not a
charge''; $\pm0.2$~eV of TiN work function is equivalent to $\pm1.12\times10^{12}\cmsq$. The 6.3~nm device-specific
value is \emph{plausible, unverified} (calibration; no independent observable).

% -------------------------------------------------------------------
\section{Effective donor density}
\label{sec:par-nd}

\paragraph{Definition and role.}
$\Ndeff(t)$ is the effective density of electrically active, ionised shallow donors in the channel (the net of
oxygen vacancies, W donors and compensating acceptors). It is \emph{not} the W concentration: the target paper
gives ``$\sim$2\,\% W'' without specifying atomic, cation or molar basis, so the composition is treated as a label and
never converted into a donor density.

\paragraph{Values used.}
\begin{equation}
  \Ndeff(t) = N_{d0}\left[1+\left(\frac{t}{t_c}\right)^{k}\right],\qquad N_{d0} = 2.5\times10^{17}\cmcu,\;
  t_c = \SI{20}{nm},\; k = 4,
  \label{eq:d1-nd}
\end{equation}
i.e.\ $2.50$ / $2.52$ / $2.97\times10^{17}\cmcu$ at 2.0 / 6.3 / 13.2~nm and $1.85\times10^{18}\cmcu$ at 31.8~nm.

\paragraph{Source.}
\prov{FITTED} law. The configuration cites as support the bulk/border trap densities 2.37 / 3.12 /
$14.6\times10^{18}\si{cm^{-3}.eV^{-1}}$ at 10 / 20 / 30~nm of a different IWO process \citep[p.~173507-6]{Kim2024}. S4 corrects
this: those numbers are gate-oxide \emph{border traps} from the carrier-number-fluctuation noise model, not donors or
IWO tail states, and the configuration ``support'' conflates the two. The credible support is S4's reading of the
same paper's turn-on voltage, which shifts linearly with $t$ by about $-2$~V per 10~nm and implies
$\Nd \approx 4.3$--$4.5\times10^{17}\cmcu$, within a factor 2 of V1. The target paper attributes electron densities
above $10^{19}\cmcu$ to undoped \InO{} and the suppression of vacancies to W \citep[p.~2]{Januar2026}; ALD \InO{}
reaches about $10^{20}\cmcu$ \citep[pp.~4--5]{Wang2022}.

\paragraph{Derivation.}
For a fully depleted film (\cref{sec:der-depletion}) with uniform donors, Poisson's equation
$d^2\psi/dz^2 = -q\Nd/\varepsilon_s$ with zero field at the free back surface ($z = t$) gives a parabolic potential
whose maximum (the conduction-band minimum) is at the back surface, $q\Nd t^2/2\varepsilon_s$ above the front value.
The gate must image the sheet $\Nd t$. Keeping the band minimum at the threshold position therefore requires
\begin{equation}
  \Delta\Vth = -\frac{q\Nd t}{\Cox} - \frac{q\Nd t^{2}}{2\varepsilon_s},
  \label{eq:d1-nd-shift}
\end{equation}
which is $-0.040$ / $-0.151$ / $-0.406$ / $-1.55$~V per $10^{18}\cmcu$ at 2 / 6.3 / 13.2 / 31.8~nm (S4; e.g.\ at 13.2~nm
$0.236 + 0.170$~V). With the V1 law the first term is only $0.009$ / $0.028$ / $0.070$~V (S1). Consequently the
donor density is \emph{unidentifiable} at 2~nm (0.036~V per $10^{18}\cmcu$, S4) and the steep $k = 4$ law matters only
above about 20~nm. The donor density could be measured by Hall or capacitance--voltage profiling on thick unpatterned films,
or from the slope of $V_{\mathrm{on}}(t)$ over a dense thickness series, which is \cref{eq:d1-nd-shift} in differential form.

\paragraph{ATLAS implementation.}
A uniform $n$-type DOPING statement on the IWO region (DOPING UNIFORM, \citealp[p.~1189]{AtlasManual}); the donor
Gaussian of DEFECTS (NGD, \citealp[p.~1171]{AtlasManual}) is zero in V1 and is the V2 candidate for a donor band at
fixed absolute energy.

\paragraph{Sensitivity evidence.}
$\Nd\times2$ (\run{0022} vs \run{0012}): $\Delta\Vthcc = -0.009$~V, fixed-current SS $+0.7\mVdec$, $\Delta\Ion = +0.7$\,\%.
A2 at 6.3~nm ($N_{d0} = 10^{16}$, \run{0032} vs \run{0014}): $+0.030$~V (analytic 0.028~V), $\SSmin$ $+6.5\mVdec$,
$\Ion$ $-1.8$\,\% -- donors are electrostatically irrelevant at 6.3~nm and reduced donors are \emph{rejected} as the
cause of the 6.3~nm miss (S1, S4). A thickness-independent donor density fitted to the 6.3$\rightarrow$13.2~nm drop
needs $2.2\times10^{18}\cmcu$ and then predicts $-0.25$~V from 2 to 6.3~nm against $-0.02$~V measured (S4). With the
V1 law, 31.8~nm gives $\Vth = -2.71$~V and $\Id(-3\,\mathrm{V}) = 1.2\times10^{-12}\Aum$, i.e.\ not always-on (S1 1-D);
its always-on state needs $\Nd \gtrsim 3\times10^{18}\cmcu$ or a back donor sheet (\cref{sec:par-backq}).

\paragraph{Status.}
\prov{FITTED}. Low $\Ndeff \approx 2.5\times10^{17}\cmcu$: \emph{plausible, unverified} (correlation with Kim2024);
a donor \emph{step} ($\approx2.5\times10^{17}$ for $t\le6.3$~nm, $\approx10^{18}$ at 13.2~nm, $\gtrsim2\times10^{18}$ at
31.8~nm) is \emph{plausible, unverified} and degenerate with $\dEc$ (S4).

% -------------------------------------------------------------------
\section{Back-surface charge and its lever arm}
\label{sec:par-backq}

\paragraph{Definition and role.}
The top surface of the channel is exposed to air. Adsorbates (O$_2$, H$_2$O) can form acceptor-like (negative) or
donor-like (positive) surface states. A sheet $\Qb$ at the back surface acts on the threshold through a larger lever
arm than a front charge, because its field crosses the film as well as the dielectric.

\paragraph{Values used.}
$\Qb = 0$ for 2.0, 6.3 and 13.2~nm in the calibrated decks. Hypothesis tests at 6.3~nm: $-1.64\times10^{12}\cmsq$
(A3, \run{0030}) and $-1.0\times10^{12}\cmsq$ within the combination A5 (\run{0035}). At 31.8~nm a donor-like back
sheet was a V0 device-specific hypothesis (configuration: centre $\Ec-0.02$~eV, width 0.05~eV, amplitude zero in V1).

\paragraph{Source.}
\prov{ASSUMED} zero; the $-1.64\times10^{12}\cmsq$ magnitude is the charge that closes the 6.3~nm $\Vthcc$ miss and
has no independent source (``adsorbate scale'' is an assumption, S1).

\paragraph{Derivation: the lever arm.}
A sheet $\Qb$ is assumed at $z = t$ and a front sheet $Q_f$ at $z = 0$, with a fully depleted film free of other charge.
The gate images $-(\Qb+Q_f)$; the field in the film between the two sheets is $q\Qb/\varepsilon_s$. Holding the
potential at the back surface fixed, a back sheet shifts the gate voltage by
\begin{equation}
  \Delta V_{\mathrm{b}} = -q\Qb\left(\frac{1}{\Cox} + \frac{t}{\varepsilon_s}\right),
  \label{eq:backq-lever}
\end{equation}
while a front sheet gives only $-q Q_f/\Cox$. The ratio $1 + \Cox t/\varepsilon_s$ is 1.22 / 1.68 / 2.44 at
2 / 6.3 / 13.2~nm (analytic back lever 0.218 / 0.301 / 0.436~V per $10^{12}\cmsq$ vs 0.179 front). \Cref{eq:backq-lever}
holds for the potential \emph{at the back surface}. The threshold current, however, flows where the electrons are:
at threshold the electron centroid is 0.8 / 2.1 / 5.7~nm from the front (S1), and the filled tail screens the back
field over the trap-screening length $L_t = 1.7$ / 3.2 / 7~nm (\cref{sec:der-depletion}). The realised lever is
therefore close to the front value: in the 1-D model with the full DOS it is 0.185 / 0.172 / 0.215~V per $10^{12}$ on
$\Vthcc$ (0.195 / 0.220 / 0.261 without the tail) and only 0.168 / 0.077 / 0.064 on $\Vthlin$, where the front
accumulation layer dominates; the back charge also lowers $\SScc$ by 11 / 38 / 47$\mVdec$ (S1).

\paragraph{ATLAS implementation.}
For A3/A5 a fixed sheet charge on the exposed IWO/air surface was added (campaign~A configuration
\file{config/campaign_A_6p3_hypotheses.json}; statement in \cref{ch:app-decks}). The V0 donor-like back sheet is a
regularised DEFECTS donor Gaussian (NGD, EGD; \citealp[pp.~1170--1171]{AtlasManual}) in a 0.25~nm IWO layer, analogous to
\cref{eq:dit-regularization}.

\paragraph{Sensitivity evidence.}
A3 (\run{0030}): $-1.64\times10^{12}\cmsq$ at the back gives $+0.278$~V against $+0.293$~V for the same charge at the
front (ratio 0.95, S1): the thick-film lever of \cref{eq:backq-lever} is \emph{not} realised at threshold. A3 reproduces
$\Vthcc$ (error $-0.017$~V) but moves $\SScc$ by $-64.1\mVdec$ and $\Vthlin$ by $-0.47$~V away from the data; the gap
$\Vthlin-\Vthcc$ becomes 0.72~V vs 1.176~V measured. A5 (\run{0035}) gives $\Vthcc$ $-0.079$~V and $\SScc$
$-45.9\mVdec$. The same charge at 2 or 13.2~nm would shift those films by 0.30 or 0.35~V, so it cannot be specific
to 6.3~nm without a sample-specific cause. At 31.8~nm the always-on state would need a back donor sheet
$\gtrsim4\times10^{12}\cmsq$ (0.80~V per $10^{12}$ at 31.8~nm; S4).

\paragraph{Status.}
\prov{ASSUMED} zero. Back-surface negative charge as the sole cause of the 6.3~nm miss: \emph{rejected}
(model-conditional, \evPOST); back-referenced lever arm at threshold: \emph{rejected} (0.96--1.2$\times$ front, not
1.7--2.4$\times$). The air-side charge of the real devices is \emph{not determined} (measurement ambient not recorded).

% -------------------------------------------------------------------
\section{Gate work function}
\label{sec:par-gatewf}

\paragraph{Definition and role.}
$\phi_M$ of the TiN gate. Together with $\chiIWO$ it sets the flat-band offset $\phi_M-\chiIWO$, which enters
$\Vthcc$ one-to-one; with the bulk affinity it contributes $+0.400$~V to every film's budget (S1).

\paragraph{Values used.}
$\phi_M = 4.70$~eV, all thicknesses.

\paragraph{Source.}
\prov{ASSUMED}, centre of a TiN literature range 4.5--4.9~eV quoted in the configuration (TiN work functions depend on
stoichiometry, deposition and the dielectric underneath); no measurement on the target stack.

\paragraph{Derivation (how it could be measured).}
The terraced-oxide (multi-EOT) flat-band method of \cref{sec:par-qf} gives $\phi_M-\chi$ as intercept; Kelvin-probe or
UPS on the TiN gives $\phi_M$ directly (with the caveat of interface dipoles). On the transfer curve it is exactly
degenerate with $\chi$ and $\Qf$: $\pm0.2$~eV $\equiv \pm1.12\times10^{12}\cmsq$ (\cref{eq:d1-qcox}; REVIEW).

\paragraph{ATLAS implementation.}
CONTACT NAME=gate WORKFUN=4.7 (the CONTACT statement specifies the physical attributes of an electrode; without it an
electrode is charge-neutral/Ohmic, \citealp[p.~1145]{AtlasManual}; WORKFUN, default 0,
\citealp[p.~1148]{AtlasManual}). The Pd source/drain are left Ohmic (\cref{sec:par-contacts}).

\paragraph{Sensitivity evidence.}
$\phi_M + 0.1$~eV: $\Delta\Vthcc = +0.100$~V at 2~nm (\run{0048} vs \run{0038}) and 13.2~nm (\run{0049} vs \run{0040}),
exact rigid shift (max $|\Delta\log_{10}\Id| = 1.6\times10^{-5}$ and $4.1\times10^{-6}$), SS unchanged ($0.0\mVdec$),
$\Delta\Ion = -6.8$ / $-5.0$\,\% at fixed $\Vg = 3$~V (S6, REVIEW).

\paragraph{Status.}
\prov{ASSUMED}; held fixed because it is degenerate with $\Qf$ and $\chi$ (\evNUM).

% -------------------------------------------------------------------
\section{Capture cross sections and recombination lifetimes}
\label{sec:par-capture}

\paragraph{Definition and role.}
Capture cross sections $\sigma_n$, $\sigma_p$ of the defect distributions and the Shockley--Read--Hall (SRH)
lifetimes of the background recombination set the \emph{rates} of trapping and generation. In a DC transfer
simulation they would matter only if trap occupancy departed from equilibrium with the electron quasi-Fermi level,
or if generation current were significant.

\paragraph{Values used.}
Tail capture cross section $10^{-15}\si{cm^2}$; background SRH lifetime $\tau = 10^{-6}$~s; both \prov{ASSUMED}, all
thicknesses (configuration). Other cross sections are those listed in the decks (\cref{ch:app-decks}); the ATLAS
defaults are, e.g., SIGTAE $10^{-16}$ and SIGTAH $10^{-14}\si{cm^2}$ \citep[p.~1169]{AtlasManual}.

\paragraph{Derivation: why DC results are independent of $\sigma$.}
For a trap level $E$ with capture coefficients $c_n = \sigma_n v_{\mathrm{th}}$, $c_p$ and emission rates
$e_n = c_n n_1$, $e_p = c_p p_1$ (detailed balance, $n_1 = \Nc e^{(E-\Ec)/\kT}$), the steady-state occupancy is
\begin{equation}
  f = \frac{c_n n + e_p}{c_n n + e_n + c_p p + e_p}
    \;\xrightarrow{\;p,\,p_1\,\rightarrow\,0\;}\;
    \frac{n}{n+n_1} = \frac{1}{1+e^{(E-E_{Fn})/\kT}} .
  \label{eq:d1-srh-occ}
\end{equation}
In the electrons-only, wide-gap $n$-type channel the hole terms vanish, $\sigma_n$ cancels, and the occupancy is the
Fermi function of the electron quasi-Fermi level. The cross sections therefore set only time constants,
$\tau_{\mathrm{trap}} = [\sigma_n v_{\mathrm{th}}(n+n_1)]^{-1}$, which matter for hysteresis, frequency dispersion and
transients -- none of which is simulated. Likewise, thermal SRH generation $q n_i t/\tau$ with $n_i\lesssim10^{-7}\cmcu$
is negligible: as an origin of the measured off-state floors it falls 13--17 decades short (Astra audit, quoted by S4).
The cross sections could be measured by transient (pulsed-$\Vg$) current recovery, deep-level transient spectroscopy,
frequency-dependent C--V or the time dependence of bias-stress shifts.

\paragraph{ATLAS implementation.}
DEFECTS SIGTAE/SIGTAH (tail) and SIGGAE/SIGGAH (Gaussian) \citep[pp.~1168--1169]{AtlasManual}; MATERIAL TAUN0
\citep[p.~1349]{AtlasManual} with SRH on the MODELS statement.

\paragraph{Sensitivity evidence.}
None run: by \cref{eq:d1-srh-occ} the DC results are invariant (\evTHEORY). The bias-stress literature for the same
material family (a threshold shift of 0.72--1.2~V after 1200~s at 6~V, as summarised by S1/S4 from the target paper)
shows that slow trapping exists; extrapolated to one $-3\ldots+3$~V sweep it gives 0.003--0.24~V, a \emph{not
determined} confounder of the 6.3~nm offset (S4).

\paragraph{Status.}
\prov{ASSUMED}; not identifiable from DC data.

% -------------------------------------------------------------------
\section{Threshold charge budget}
\label{sec:der-vth-budget}

\paragraph{Purpose.}
This derivation decomposes the simulated constant-current threshold $\Vthcc$ (at $10^{-9}\Aum$, $\Vd = 0.7$~V;
\cref{sec:der-metrics}) into additive electrostatic terms, so that every input of this chapter has a quantified weight
and the thickness dependence can be attributed term by term.

\paragraph{Derivation.}
At threshold the film is fully depleted and nearly flat-band inside (band bending across the film 0.022 / 0.045 /
0.016~eV in the ATLAS profiles of \run{0012} / \run{0015} / \run{0013}, S1). Gauss's law across the dielectric, with
the gate Fermi level at $-q\Vg$ relative to the source electron quasi-Fermi level, gives
\begin{equation}
  \Vthcc = \underbrace{\frac{\phi_M-\chi^{\mathrm{bulk}}}{q}}_{\text{alignment}}
         + \underbrace{\frac{\dEc}{q}}_{\text{confinement}}
         + \underbrace{\frac{(E_{Fn}-\Ec)_{\mathrm{cc}}}{q}}_{\text{sets }n_s}
         - \frac{q\Qf}{\Cox}
         - \frac{q\Nd t}{\Cox}
         + \frac{q\,(n_t + N_{\mathrm{it}}^{-} + n_{\mathrm{GA}}^{-} + n_s)}{\Cox},
  \label{eq:d1-vth-budget}
\end{equation}
where $n_t$, $N_{\mathrm{it}}^{-}$, $n_{\mathrm{GA}}^{-}$ are the ionised tail, interface and deep-Gaussian sheets
(\cref{eq:d1-nt-occ,eq:d1-gauss,eq:dit-regularization}), $n_s$ the free sheet at the source, and
$(E_{Fn}-\Ec)_{\mathrm{cc}}$ the Fermi-level position at which the constant-current criterion is met. The last term
depends on mobility: a film with higher $\muband$ reaches $10^{-9}\Aum$ at a lower $n_s$ and hence a lower $\EF$, which
also empties part of the tail. S1 evaluates \cref{eq:d1-vth-budget} with a 1-D Poisson surrogate that matches ATLAS
$\Vthcc$ within 3~mV (runs \run{0015}, \run{0030}, \run{0013}); its closure is exact to $10^{-6}$~V.

\begin{table}[tbp]
  \centering\footnotesize\setlength{\tabcolsep}{3pt}
  \caption{Threshold charge budget at $\Vthcc$ (V), S1 1-D surrogate (\file{s1_poisson1d.out.txt}); \cref{fig:vth_budget}
  shows the same terms graphically. Sheet densities in parentheses (cm$^{-2}$). Columns: 2.0~nm \run{0012}; 6.3~nm
  validation \run{0014}; 6.3~nm tuned \run{0015}; 13.2~nm \run{0013}.}
  \label{tab:d1-budget}
  \begin{tabular}{lcccc l}
    \toprule
    term & 2.0 nm & 6.3 nm val. & 6.3 nm tuned & 13.2 nm & status \\
    \midrule
    $(\phi_M-\chi^{\mathrm{bulk}})/q$ & $+0.400$ & $+0.400$ & $+0.400$ & $+0.400$ & assumed \\
    $\dEc/q$ & $+0.353$ & $+0.072$ & $+0.072$ & $+0.026$ & DFT proxy \\
    $(E_{Fn}-\Ec)_{\mathrm{cc}}/q$ & $-0.038$ & $-0.039$ & $-0.037$ & $-0.111$ & model \\
    $-q\Qf/\Cox$ & $-0.310$ & $-0.310$ & $-0.016$ & $-0.310$ & fitted \\
    $-q\Nd t/\Cox$ & $-0.009$ & $-0.028$ & $-0.028$ & $-0.070$ & fitted law \\
    tail $q n_t/\Cox$ & \makecell{$+0.256$\\($1.43\times10^{12}$)} & \makecell{$+0.226$\\($1.27\times10^{12}$)} & \makecell{$+0.231$\\($1.29\times10^{12}$)} & \makecell{$+0.098$\\($5.5\times10^{11}$)} & fitted \\
    interface traps & $+0.012$ & $+0.012$ & $+0.012$ & $+0.012$ & assumed \\
    free electrons & \makecell{$+0.014$\\($7.6\times10^{10}$)} & \makecell{$+0.016$\\($9.1\times10^{10}$)} & $+0.017$ & \makecell{$+0.005$\\($2.6\times10^{10}$)} & model \\
    deep Gaussian & $0.000$ & $+0.001$ & $+0.001$ & $+0.003$ & assumed \\
    \midrule
    $\Vthcc$ 1-D / ATLAS / measured & 0.679 / 0.676 / 0.662 & 0.351 / 0.350 / 0.644 & 0.652 / 0.651 / 0.644 & 0.052 / 0.053 / 0.083 & \\
    \bottomrule
  \end{tabular}
\end{table}

\paragraph{Attribution of the thickness dependence.}
Differencing \cref{tab:d1-budget} between films gives the model's thickness steps (S1): 2$\rightarrow$13.2~nm,
$-0.627$~V in total, of which confinement $-0.327$ (52\,\%), tail $-0.158$ (25\,\%), $(E_{Fn}-\Ec)$ $-0.073$ (12\,\%) and
donors $-0.061$ (10\,\%), against $-0.579$~V measured; 2$\rightarrow$6.3~nm, $-0.328$~V, of which confinement $-0.281$
(86\,\%), tail $-0.030$, $(E_{Fn}-\Ec)$ $-0.001$ and donors $-0.019$, against $-0.018$~V measured. Hence the V1 6.3~nm
miss (\run{0014}: $-0.294$~V) is essentially the confinement difference $\dEc(2)-\dEc(6.3)$ (\emph{demonstrated within
the model}, S1). Only two analytic terms are fixed by laws rather than fitted: $q\Nd t/\Cox$ and $\dEc$. The
constant-current definition itself contributes: if the 13.2~nm film had the 2~nm mobility its $\Vthcc$ would be
0.108--0.115~V higher (analytic on data), so 0.11--0.14~V of the measured $-0.58$~V step is a consequence of the
threshold definition, not of charge (\emph{strongly supported}, S1; naive estimate $\SScc\log_{10}(\mu$ ratio$) = 0.10$~V
at $160\mVdec$). One-at-a-time removal in the 1-D model confirms the weights: confinement ($\dEc$ and $\mstar$) 0.315 /
0.064 / 0.023~V, tail 0.221 / 0.215 / 0.112~V, $\mstar$ via $\Nc$ only $-0.039$ / $-0.008$ / $-0.002$~V, $\Nd$ $-0.009$ /
$-0.031$ / $-0.094$~V, $\Dit$ 0.012 / 0.012 / 0.011~V (2 / 6.3 tuned / 13.2~nm), consistent with the ATLAS differences
\run{0012}$-$\run{0016} (0.315~V), \run{0013}$-$\run{0017} (0.023~V), \run{0022} and \run{0025}.

\paragraph{Degenerate rigid shifts.}
$\phi_M$, $\chi$, $\Qf$, $\dEc$ (to within the $\Nc$ term), $\Dit$ charge, deep-Gaussian charge and $\Cox$ errors all enter
\cref{eq:d1-vth-budget} as rigid offsets at first order; on one $\Id$--$\Vg$ curve only their weighted sum is identified
(at 6.3~nm about $1.64\times10^{12}\cmsq \equiv 0.29$~V; S1). This is the formal statement of the non-identifiability of
the confinement shift.

\paragraph{Band diagrams.}
\Cref{fig:band_2p0_eq,fig:band_2p0_vth,fig:band_2p0_on,fig:band_13p2_on} show the ATLAS band diagrams of the calibrated
decks, on which the budget is based: the film is nearly flat inside at threshold, so the budget terms are additive, and
the gate field bends the bands only in the on-state.

\begin{figure}[tbp]
  \centering
  \begin{subfigure}[t]{0.49\linewidth}
    \includegraphics[width=\linewidth]{band_2p0_eq}
    \caption{2.0~nm, equilibrium.}
    \label{fig:band_2p0_eq}
  \end{subfigure}\hfill
  \begin{subfigure}[t]{0.49\linewidth}
    \includegraphics[width=\linewidth]{band_2p0_vth}
    \caption{2.0~nm, at threshold.}
    \label{fig:band_2p0_vth}
  \end{subfigure}\\[4pt]
  \begin{subfigure}[t]{0.49\linewidth}
    \includegraphics[width=\linewidth]{band_2p0_on}
    \caption{2.0~nm, on-state.}
    \label{fig:band_2p0_on}
  \end{subfigure}\hfill
  \begin{subfigure}[t]{0.49\linewidth}
    \includegraphics[width=\linewidth]{band_13p2_on}
    \caption{13.2~nm, on-state.}
    \label{fig:band_13p2_on}
  \end{subfigure}
  \caption{Energy-band diagrams of the calibrated decks \run{0012} (2.0~nm, panels a--c) and \run{0013} (13.2~nm, panel d),
  as exported by the package plotting scripts (\file{plots/2p0/band_diagram_{eq,vth,on}.png},
  \file{plots/13p2/band_diagram_on.png}). The panels show the conduction-band edge of the 2~nm film raised by confinement;
  the almost flat bands inside the film at threshold (ATLAS band bending across the film 0.022~eV at 2~nm and 0.016~eV
  at 13.2~nm, S1), which justifies the additive budget of \cref{tab:d1-budget}; and the front-interface band bending in
  the on-state.}
\end{figure}

\paragraph{Status.}
\evSURR{} checked against \evNUM{} (3~mV). The budget explains the model, not the device: the measured 2 $\approx$ 6.3
$\gg$ 13.2~nm pattern of $\Vthcc$ is not reproduced by any single-material law (\cref{ch:mechanisms}).

% -------------------------------------------------------------------
\section{Subthreshold slope relations and trap-DOS inversion}
\label{sec:der-ss}

\paragraph{Charge-control derivation.}
Below threshold the current is proportional to the free sheet density at the source, $\Id\propto n_s\propto
e^{u/\kT}$ with $u = E_{Fn}-\Ec$ (non-degenerate). Differentiating the Gauss law
$\Cox(\Vg - V_{\mathrm{FB}} - u/q) = q(n_s + n_t + N_{\mathrm{it}}^{-} + \ldots)$ gives
$\Cox\,d\Vg = (\Cox/q)\,du + q\,dn_s + q\,dn_t + q\,dN_{\mathrm{it}}^{-}$, so
\begin{equation}
  SS \equiv \frac{d\Vg}{d\log_{10}\Id}
     = \ln 10\,\frac{\kT}{q}\left[1 + \frac{C_{\mathrm{it}} + C_{t} + C_{\mathrm{free}}}{\Cox}\right],\qquad
  C_{\mathrm{it}} = q^{2}\Dit,\quad C_t = q^{2}t\,\frac{\partial n_t^{\mathrm{vol}}}{\partial u},\quad
  C_{\mathrm{free}} = q^{2}\frac{\partial n_s}{\partial u} = \frac{q^{2}n_s}{\kT}.
  \label{eq:ss-dit}
\end{equation}
For the exponential tail, $C_t = q^2 t\,n_t^{\mathrm{vol}}/\WTA$ (\cref{eq:d1-nt-occ}), so the tail contribution grows
exponentially as $\EF$ approaches $\Ec$; the free-carrier term takes over once $n_s/\kT$ exceeds $\Cox/q^2$, the
crossover ``$C_{\mathrm{free}} = \Cox$''. The textbook inversion of \cref{eq:ss-dit},
\begin{equation}
  D_{\mathrm{tot}} = \frac{\Cox}{q^{2}}\left(\frac{SS}{\ln 10\,\kT/q} - 1\right),
  \label{eq:d1-ss-naive}
\end{equation}
attributes the whole excess slope to traps. It is the method used by the literature values of
\cref{sec:par-dit} (e.g.\ the \HfO/\InO{} $\Dit$ from an SS of 63.5$\mVdec$, \citealp[p.~21539]{Lin2022}).

\paragraph{Why the naive inversion fails here.}
Applied to the measured slopes, \cref{eq:d1-ss-naive} gives $2.35$ / $5.18$ / $6.69\times10^{12}\si{cm^{-2}.eV^{-1}}$
from $\SSmin$ and $1.98$ / $2.22$ / $0.94\times10^{13}$ from $\SScc$ (2 / 6.3 / 13.2~nm; S4). Two confounders make these
numbers unreliable (S4): (i) \emph{mobility}: the crossover $C_{\mathrm{free}} = \Cox$ sits at $3.2\times10^{-10}$ /
$4.0\times10^{-10}$ / $1.8\times10^{-9}\Aum$ for $\mu = 10.7$ / 13.3 / 59$\cmVs$, so the $10^{-10}$--$10^{-8}$ window
straddles it; a trap-free 2~nm film already gives $\SScc = 122\mVdec$ at $\mu = 13.3$ against 75 at $\mu = 59$, i.e.\
about 60$\mVdec$ of the thin films' $\SScc$ (about 15 at 13.2~nm) is mobility, not traps. (ii) \emph{energy window}:
the low 13.2~nm $\SScc$ reflects a window about $\kT\ln4.5 \approx 39$~meV deeper in energy, not fewer traps. The naive
formula therefore ranks 13.2~nm lowest ($9.4\times10^{12}$ vs $2.0\times10^{13}$), whereas the proper inversion ranks it
highest. The floor-dependent $\SSmin$ window (\cref{sec:der-metrics}) adds a further artefact: it reaches $u\approx-0.21$ / $-0.17$ / $-0.13$~eV
for 2 / 6.3 / 13.2~nm because the off-floors rise $\times1432$ (\cref{fig:off_floors}).

\paragraph{Charge-sheet trap-DOS inversion.}
S4 removes both confounders with an exact identity. In the gradual-channel approximation
$\Id = (W/L)\,q\mu\int_0^{\Vd} n_s(\Vg - V)\,dV$, hence $\gm = (W/L)\,q\mu\,[n_s(\Vg) - n_s(\Vg-\Vd)]$ and, by telescoping,
\begin{equation}
  n_{s0}(\Vg) = \frac{L}{q\mu W}\sum_{k\ge0}\gm(\Vg - k\Vd),
  \label{eq:d1-ns-sum}
\end{equation}
which removes the drain term. With $u$ obtained from $n_{s0}$ through $\Nc$ and Fermi--Dirac statistics, differentiating
the Gauss law gives the trapped DOS
\begin{equation}
  D_{\mathrm{trap}}(u) = \frac{\Cox}{q^{2}}\left(q\,\frac{d\Vg}{du} - 1\right) - \frac{dn_f}{du},
  \label{eq:d1-dtrap}
\end{equation}
in which $V_{\mathrm{FB}}$ cancels. The method was validated on ATLAS curves before use (measured/ATLAS DOS ratio 0.81--1.17
for $u = -0.125$ to $-0.025$~eV at 2 and 13.2~nm) and requires the mobility, which shifts the energy axis by $\kT\ln(\mu$
ratio$)$, $\pm10$--15~meV between the V1 and S2 mobilities (S4). Results: near $\Ec$ the sheet DOS grows sub-linearly with
thickness, $D(\Ec-0.1\,\mathrm{eV}) = 1.05$ / 1.67 / $2.16\times10^{13}\si{cm^{-2}.eV^{-1}}$ ($\propto t^{0.38\pm0.15}$), a
surface\,+\,bulk signature (\cref{sec:par-tail}); V1 underestimates the states at $\Ec-0.15$ to $-0.2$~eV by $\times2.2$--2.8
(\cref{sec:par-deep}); and the 6.3~nm film carries more states in the last $\sim$50~meV below $\Ec$ (ratios 1.25 / 1.46 /
1.72 at $-0.05$ / $-0.025$ / 0~eV relative to \run{0015}), which fill by $\Vthlin$ and then saturate -- a later, steeper
turn-on, as measured (S4; \cref{ch:6p3}).

\begin{figure}[tbp]
  \centering
  \includegraphics[width=\linewidth]{trap_spectroscopy}
  \caption{Trap-DOS spectroscopy by the charge-sheet inversion (\cref{eq:d1-ns-sum,eq:d1-dtrap}) of S4 (copied from
  \file{scratch/S4/s4_trap_spectroscopy.png}): trapped-state density vs $u = E_{Fn}-\Ec$ from the measured curves and from
  the simulated curves of the calibrated decks. The figure shows the agreement of measured and ATLAS DOS within 0.81--1.17 for
  $u = -0.125$ to $-0.025$~eV (2 and 13.2~nm); the sub-linear growth of the near-edge sheet DOS with thickness; the
  underestimate of the deeper states by V1.}
  \label{fig:trap_spectroscopy}
\end{figure}

\paragraph{Sensitivity evidence.}
Slope-changing inputs at 2~nm (vs \run{0012}; fixed-current SS in $10^{-10}$--$10^{-9}\Aum$): $\Nt\times1.5$ $+49.0$,
$\WTA+5$~meV $+4.6$, $\Dit\times3$ $+0.1$, $\mstar\times1.3$ $-18.0\mVdec$ (\run{0050} vs \run{0038}); rigid inputs
($\chi$, $\phi_M$, $\Qf$) change it by $0.0$. The surrogate reproduces ATLAS $\SScc$ as 294 / 303 / 165 vs
289 / 289 / 193$\mVdec$ (S4 part C).

\paragraph{Status.}
\evTHEORY{} (\cref{eq:ss-dit,eq:d1-ns-sum,eq:d1-dtrap}); inversion \evSURR{} validated against \evNUM. $\SScc(t)$ read as a trap
trend with \cref{eq:d1-ss-naive}: \emph{rejected} (S4). $\SSmin$ is retired as a metric (\cref{sec:der-metrics}).

% -------------------------------------------------------------------
\section{Debye length, depletion width and full depletion}
\label{sec:der-depletion}

\paragraph{Debye and trap-screening lengths.}
Linearising Poisson's equation about a uniform electron density $n$ gives the screening length
\begin{equation}
  L_D = \sqrt{\frac{\varepsilon_s\,\kT}{q^{2}n}},\qquad
  L_t = \sqrt{\frac{\varepsilon_s\,\WTA}{q^{2}n_t^{\mathrm{vol}}}},
  \label{eq:d1-debye}
\end{equation}
the second for screening by the exponential tail ($\partial n_t/\partial u = n_t/\WTA$). With $\varepsilon_s = 9.3$:
$L_D = 7.3$ / 6.7 / 3.7 / 2.1~nm at $n = 2.5\times10^{17}$ / $3\times10^{17}$ / $10^{18}$ / $3\times10^{18}\cmcu$, and
$L_t = 1.7$ / 3.2 / 7~nm at the threshold trapped densities of 2 / 6.3 / 13.2~nm ($7.2\times10^{18}$ / $2.0\times10^{18}$ /
$4.2\times10^{17}\cmcu$) (S1). The regularisation layer of \cref{eq:dit-regularization} (0.25~nm) is well below both.

\paragraph{Depletion width and full depletion.}
A surface sheet $Q_s$ is compensated by depleting donors over $W = Q_s/\Nd$; the maximum depletion width for a band
bending $\psi$ is $W_{\max} = \sqrt{2\varepsilon_s\psi/(q\Nd)}$. At the V1 $\Nd$ ($2.5$--$3\times10^{17}\cmcu$) a sheet of
$10^{12}\cmsq$ depletes 40~nm, so every film $\le31.8$~nm is fully depleted (S1). A neutral (ungated) layer between 6.3
and 13.2~nm would require $7.6\times10^{17} < \Nd < 1.6\times10^{18}\cmcu$ (with $Q_s\approx10^{12}$). The data exclude
it at 13.2~nm: the off-band current $6.6\times10^{-12}\Aum$ bounds any ungated free sheet to
$n_s = I L/(q\mu\Vd W) \le 2\times10^{7}\cmsq$ ($\mu = 59\cmVs$), whereas a neutral sheet of only $10^{10}\cmsq$ would
carry $3.3\times10^{-9}\Aum$ (S1). The always-on 31.8~nm film needs at least $5.6$--$9.2\times10^{11}\cmsq$ ungated,
i.e.\ uniform $\Nd\ge3\times10^{18}\cmcu$, which would put 13.2~nm at $-1.0$~V; the 13.2$\rightarrow$31.8~nm transition
therefore requires a $\ge10\times$ change of effective donor density or a surface donor layer -- a material step, not
dimensional electrostatics (S1).

\paragraph{Where the electrons are.}
Subthreshold conduction is volumetric: the electron centroid at threshold is 0.8 / 2.1 / 5.7~nm from the front
(2 / 6.3 / 13.2~nm). In the on-state at 3~V all films carry a 0.7--1.1~nm front accumulation layer (ATLAS profiles; S1),
which is where the field quantization of \cref{sec:der-confinement} acts. \Cref{fig:edens_2p0_vs_vg,fig:edens_13p2_vs_vg}
show the evolution of the electron density with gate voltage in the calibrated decks.

\begin{figure}[tbp]
  \centering
  \begin{subfigure}[t]{0.49\linewidth}
    \includegraphics[width=\linewidth]{edens_2p0_vs_vg}
    \caption{2.0~nm (\run{0012}).}
    \label{fig:edens_2p0_vs_vg}
  \end{subfigure}\hfill
  \begin{subfigure}[t]{0.49\linewidth}
    \includegraphics[width=\linewidth]{edens_13p2_vs_vg}
    \caption{13.2~nm (\run{0013}).}
    \label{fig:edens_13p2_vs_vg}
  \end{subfigure}
  \caption{Electron density vs gate voltage in the calibrated decks, as exported by the package plotting scripts
  (\file{plots/2p0/electron_density_vs_vg.png}, \file{plots/13p2/electron_density_vs_vg.png}). The
  thin film fills uniformly (volumetric conduction, centroid $\approx t/2$), whereas the 13.2~nm film develops a front
  accumulation layer of about 1~nm at high $\Vg$ (S1).}
\end{figure}

\paragraph{Status.}
\evTHEORY{} with \evSURR/\evNUM{} values. Full depletion of all films $\le31.8$~nm at the V1 $\Nd$: \emph{demonstrated within
the model}; a neutral bulk layer as the 6.3$\rightarrow$13.2~nm step: \emph{rejected} (S1).

% -------------------------------------------------------------------
\section{Summary}
\label{sec:d1-summary}

\begin{table}[tbp]
  \centering\footnotesize
  \caption{Electrostatic inputs of the V1 model: values, status and controlling runs.}
  \label{tab:d1-summary}
  \begin{tabularx}{\linewidth}{l l X l}
    \toprule
    input & 2.0 / 6.3 / 13.2 nm & sensitivity (runs) & status \\
    \midrule
    $\Cox$, EOT & \SI{8.955e-7}{F/cm^2}, 3.86~nm & analytic, $q/\Cox = 0.179$~V per $10^{12}$ & \prov{CALCULATED} \\
    $\epsIWO$ & 9.3 & $<1$\,\% at 2~nm (\run{0026}) & \prov{MEASURED\_TARGET\_DEVICE} \\
    $\Eg$ & 3.403 / 3.122 / 3.076~eV & none (holes not solved) & \prov{LITERATURE\_IWO} + proxy \\
    $\chiIWO$ & 3.947 / 4.228 / 4.274~eV & $\mp0.100$~V rigid (\run{0020}/\run{0021}) & \prov{LITERATURE\_IWO} + proxy \\
    $\dEc$ & 0.353 / 0.072 / 0.026~eV & $-0.315$ / $-0.023$~V if removed (\run{0016}, \run{0017}) & \prov{DFT\_PURE\_IN2O3\_PROXY} \\
    $\mstar$ & 0.257 / 0.218 / 0.211~$m_0$ & $\times1.3$: $-0.044$ / $-0.029$~V (\run{0050}/\run{0051}) & proxy \\
    $\Nc$ & 3.28 / 2.55 / $2.44\times10^{18}\cmcu$ & via $\mstar$ & \prov{CALCULATED} \\
    $\Nv$ & $10^{19}\cmcu$ & none & \prov{ASSUMED} \\
    $\Nt$, $\WTA$ & 2.0 / 0.85 / $0.49\times10^{19}\cmcu$; 40~meV & $\times1.5$: $+0.082$~V, $-25$\,\% $\Ion$ (\run{0023}); $+5$~meV: $+0.020$~V (\run{0024}) & \prov{FITTED} \\
    deep Gaussian & $5\times10^{16}$, $\Ec-0.6$, 0.15~eV & $\le0.003$~V (analytic) & \prov{ASSUMED} \\
    $\Dit$ & $3\times10^{11}$ at $\Ec-0.3$, 0.12~eV & $\times3$: $+0.025$~V (\run{0025}); $\times5$: $+0.051$~V (\run{0033}) & \prov{ASSUMED} \\
    $\Qf$ & $1.73\times10^{12}$ (6.3 tuned $8.7\times10^{10}$) & exact rigid $-0.309$~V (\run{0001}$\rightarrow$\run{0003}) & \prov{FITTED} \\
    $\Ndeff$ & 2.50 / 2.52 / $2.97\times10^{17}\cmcu$ & $\times2$: $-0.009$~V (\run{0022}); A2 $+0.030$~V (\run{0032}) & \prov{FITTED} \\
    $\Qb$ & 0 & A3 $-1.64\times10^{12}$: $+0.278$~V, SS $-64$ (\run{0030}) & \prov{ASSUMED} \\
    $\phi_M$ & 4.70~eV & $+0.1$~eV: $+0.100$~V rigid (\run{0048}/\run{0049}) & \prov{ASSUMED} \\
    $\sigma$, $\tau_{\mathrm{SRH}}$ & $10^{-15}\si{cm^2}$, $10^{-6}$~s & none (DC-invariant) & \prov{ASSUMED} \\
    \bottomrule
  \end{tabularx}
\end{table}

\Cref{tab:d1-summary} collects the inputs. Only the gate stack (geometry and $\varepsilon_{\mathrm{HfO_2}}$, hence most
of $\Cox$) and $\epsIWO$ are measured on the target process; the band-structure inputs are literature values of other IWO devices or pure-\InO{} DFT proxies; the
charge-and-trap inputs are fitted or assumed. The model is nonetheless well conditioned for what it is used for,
because the inputs separate into two classes with different signatures: \emph{rigid shifts} ($\phi_M$, $\chi$, $\Qf$,
$\dEc$ within the $\Nc$ term, $\Dit$ and deep-Gaussian charge, $\Cox$ errors), of which only a weighted sum is identified
per curve; and \emph{shape parameters} ($\Nt$, $\WTA$, $\mstar$ through $\Nc$, near-edge bands), which the slope and the
$\Vthlin-\Vthcc$ gap constrain.

\begin{keyresult}
The electrostatics of the V1 IWO model is fixed by a measured gate stack ($\Cox = \SI{8.955e-7}{F/cm^2}$, so
$0.179$~V per $10^{12}\cmsq$), a DFT-proxy confinement law ($\dEc = 0.353$ / 0.072 / 0.026~eV) and a fitted exponential
acceptor tail ($\Nt = 2\times10^{19}(2/t)^{0.75}\cmcu$, $\WTA = 40$~meV) with a shared fixed charge
$\Qf = 1.73\times10^{12}\cmsq$. The threshold budget (1-D surrogate within 3~mV of ATLAS) attributes 52\,\% of the model's
2$\rightarrow$13.2~nm step and 86\,\% of its 2$\rightarrow$6.3~nm step to $\dEc$, and 25\,\% of the former to the tail
(\evSURR/\evNUM). Confinement at 2~nm is theoretically expected (Schr\"odinger--Poisson 0.14--0.38~V, central 0.26--0.30~V),
but it is exactly the kind of rigid shift that one transfer curve cannot distinguish from $\approx1.7\times10^{12}\cmsq$ of
fixed charge, and the $t^{-1.38}$ law overestimates the subthreshold-equivalent shift by $\approx2\times$ at 6.3~nm. $\phi_M$, $\chi$ and $\Qf$ are one
parameter (exact rigid shifts, runs \run{0020}, \run{0021}, \run{0048}, \run{0049}); donors, permittivity, interface traps and
the deep Gaussian are electrostatically small at $t\le6.3$~nm (runs \run{0022}, \run{0026}, \run{0025}, \run{0032}); the tail
controls slope and threshold. Status: \emph{calibrated, not validated}.
\end{keyresult}

The electrostatics fixes where the Fermi level sits for a given gate voltage; how much current flows then depends on the
mobility, the contacts and the extraction of the metrics, which are derived in the next chapter.
